%!TEX root = ./../../springer-mono-format/book.tex
%!TEX root = ./../../thesis-trier.tex


\begin{flushright}
	\itshape %\setstretch{1.0}
	``It was at this point that Wigner, Hund, Heitler, and Weyl entered the picture with their `Gruppenpest': the pest of the group theory… The authors of the `Gruppenpest' wrote papers which were incomprehensible to those like me who had not studied group theory, in which they applied these theoretical results to the study of the many electron problem. The practical consequences appeared to be negligible, but everyone felt that to be in the mainstream one had to learn about it. Yet there were no good texts from which one could learn group theory. It was a frustrating experience, worthy of the name of a pest.''  \\ \normalfont
	\vspace{.2em}
	\rightskip=.8cm---John Slater (1900--1976)
\end{flushright}

The set $\mathcal S_X$ of all permutations on a set $X$
with function composition
and inversion given by $t^{-1} \colon X \to X$ 
for $t \in \mathcal S_X$ is a group, called the \emph{(full) symmetric group on $X$}.
%
%\bigskip 
%
%For a set $X$, the \emph{(full) symmetric group on $X$} is the group structure on the set~$\mathcal S_X$ of all permutations on~$X$
%with function composition as operation
%and inversion given by $t^{-1} \colon X \to X$ 
%for $t \in \mathcal S_X$.
%
%\bigskip
%
%For a set $X$, the set $\mathcal S_X$ of all permutations on $X$,
%which is a group with function composition
%and inversion given by $t^{-1} \colon X \to X$ 
%for $t \in \mathcal S_X$, called the \emph{(full) symmetric group on $X$}.
A \emph{permutation group~$G$} (over the \emph{(permutation) domain}~$X$) is a subgroup of~$\mathcal S_X$.
Here, we investigate permutation groups.

\Cref{prelim:sec:basic} introduces basic notions. In \cref{prelim:permgrp:sec:transitive}
we collect statements related to transitive permutation groups.
The relation between permutation groups and (abstract) groups
is discussed in \cref{prelim:sec:perm_grp_abstract_grp}. \Cref{prelim:sec:grp-of-units} investigates
the group of units in transformation semigroups.
Lastly, \cref{prelim:sec:graphs-groups} points out some relations between permutation groups and graphs.

\section{Basic Notions}
\label{prelim:sec:basic}

%!TEX root = ../../springer-mono-format/book.tex
%!TEX root = ../../thesis-trier.tex

\IN{symmetric group} 
\IN{full symmetric group}
\IN{permutation group}
\IN{degree of permutation group}
\IN{domain!of permutation group} 
\label{def:S_X}


Let $G \le_G \mathcal S_X$ be a permutation group.

\begin{enumerate}
	\item If $X$ is finite, we call $|X|$ the \emph{degree (of $G$)}.
	For $n\ge 0$, we set $\mathcal S_n = \mathcal S_{[n]}$. \label{def:S_n} 
	
	\label{def:block-perm-grp}
	\IN{block!in permutation group} 
	\IN{system of blocks} 
	\IN{trivial block}
	\item A non-empty subset $\Delta \subseteq X$ is a \emph{block}\footnote{To distinguish them from ordinary sets, blocks for $G$ 
		are denoted by capital Greek letters. Also, alphabets are denoted by capital
		Greek letters, but the context should always make it clear
		whether a specific capital Greek letter denotes a block or an alphabet.} if and only if, for all $g \in G$,
	the condition $\applyfunc{\Delta}{g} \cap \Delta \ne \emptyset$ implies $\applyfunc{\Delta}{g} = \Delta$.
	A collection of blocks that forms a partition of~$X$
	is called a \emph{system of blocks}. 
	Blocks are precisely the congruence classes of some $G$-congruence on $X$.
	A block is a \emph{trivial block} if it is a singleton or equals $X$.
	
	\begin{comment} 
		A family blocks $\{ X_i \mid i \in I\}$ such that $|I| \ge 2$ and t
		\begin{enumerate}
			\item 
			
			
		\end{enumerate}
		
		containing at least two blocks and such that least one block contains at least two elements
		whose union equals $X$ is called a \emph{system of imprimitivity}.
	\end{comment}






\IN{primitive permutation group}
\IN{imprimitive permutation group} 
\item A permutation group $G$ over $X$ is \emph{primitive} if the only
$G$-congruences on~$X$ are the trivial $G$-congruences.
As the blocks are precisely the $G$-congruence classes
for some $G$-congruence, 
$G$ is primitive if and only if it only admits
the trivial blocks.
If $G$ is transitive and not primitive, then we call it \emph{imprimitive}.
		
	\IN{$G$-orbit} 
	\IN{orbit|see{$G$-orbit}}
	\item The \emph{$G$-orbit} (or just \emph{orbit}) of $a \in X$ in $G$ is $\Orbit{a}{G} = \{ \applyfunc{a}{g} \mid g \in G\}$. The orbits $\Orbit{a}{G}$ for $a \in X$ form a system of blocks,
	and the union of several orbits is a block.
	\label{def:orbit}
		
	\IN{$k$-homogeneous}
	\IN{$k$-transitive}
	\IN{intransitive permutation group}
	\IN{transitive on subset} 
	\IN{transitive permutation group}
	\IN{acts transitively}
	\item Let $k \ge 1$. Then $G$ is called 
	\emph{$k$-homogeneous} (resp.~\emph{$k$-transitive})
	if for every two $k$-sets $A, B \subseteq X$ (resp.~$k$-tuples $(a_1, \ldots, a_k)$, $(b_1, \ldots, b_k) \in X^k$)
	there exists $g \in G$ such that 
	$B = \applyfunc{A}{g}$ (resp.~$(b_1, \ldots, b_k) = (g(a_1), \ldots, g(a_k))$).
	Every $k$-transitive group is $k$-homogeneous.
	The symmetric group~$\mathcal S_{X}$ is $k$-transitive 
	for each $k \ge 1$.
	A $1$-homogeneous permutation group is also called \emph{transitive}.
	Equivalently,~$G$ is transitive  
	if and only if there exists $a \in X$
	such that $X = \Orbit{a}{G}$. 
	%We also say that $G$ \emph{acts transitively} in this case. 
	For $A \subseteq X$ we also say that~$G$ is \emph{transitive on~$A$}, or \emph{acts transitively on $A$},
	if $\Orbit{a}{G} = A$ for some (and so each) $a \in A$.	
	We call $G$ \emph{intransitive} if it is not transitive.
	
	\IN{point stabilizer}
	\item Let $a \in X$. The \emph{point stabilizer subgroup (of $a$ in $G$)} is $\Stab{a}{G} = \{ g \in G \mid \applyfunc{a}{g} = a \}$. \label{def:stabilizer}
	For $g \in G$ we have 
	\begin{equation}\label{prelim:coset-point-stab}
		g\Stab{a}{G} = \{ h \in G \mid \applyfunc{a}{h} = \applyfunc{a}{g} \}.
	\end{equation}
			
	\IN{conjugate transformation}
	\label{def:conjugation}
	\item For any transformation $t \colon X \to X$ and permutation $g \colon X \to X$ 
	the \emph{conjugate of $t$ by $g$} is $t^g \colon X \to X$ given by 
	$\concat{g^{-1}}{\concat{t}{g}}$. For $a,b \in X$, we have\footnote{Note that by our convention of reading functional composition from right to left, the inverse has to occur on the right hand side in the defining equation of $t^g$ to make \cref{prelim:conj} valid. In the literature, for example in~\cite{Kurzweil04}, 
	conjugation of~$t$ by~$g$ is also often defined as $g^{-1} t g$ instead.}
	\begin{equation}\label{prelim:conj}
		\applyfunc{a}{t} = b \Leftrightarrow 
		\applyfunc{\applyfunc{a}{g}}{t^g} = \applyfunc{b}{g}.	
	\end{equation}
	This equation implies, for $a \in X$ and $g \in G$,
	\begin{equation}\label{prelim:eqn:stab}
		\Stab{\applyfunc{a}{g}}{G} = g\Stab{a}{G}g^{-1}.
	\end{equation}		
	By \cref{prelim:conj}
	every permutation $g \colon X \to X$ 
	is a graph morphism of the functional graphs of $t$
	and $t^g$ for a transformation $t \colon X \to X$.	
	
	\IN{semiregular permutation group} 
	\item We call $G$ \emph{semiregular}
	if $\Stab{a}{G}$ is trivial for each $a \in X$.
	% Equivalently, for all $a \in X$, if $\applyfunc{a}{g} = \applyfunc{a}{h}$
	% for $g,h \in G$, then $g = h$.
	This is equivalent to the condition that $\applyfunc{a}{g} = \applyfunc{a}{h}$
    for $g,h \in G$ and some $a \in X$ implies $g = h$
    since $g(a) = h(a)$ yields $h^{-1}(g(a)) = a$.
    By semiregularity, we find $h^{-1}(g) = \operatorname{id}_X$,
    giving $g = h$.
	
	\IN{regular permutation group}
	\item We call $G$ \emph{regular}  
	if it is semiregular and transitive.
	Equivalently, $G$ is transitive and there exists $a \in X$
	such that $\Stab{a}{G} = \{ \operatorname{id}_{X}\}$.
		
	\IN{alternating group}
	\label{def:A_n}
	\item The \emph{alternating group} (for a given $n$) 
	is the unique subgroup $\mathcal A_n \le_G \mathcal S_n$ such that $|\mathcal A_n| = |\mathcal S_n| / 2$. 
	See~\cite{cameron_1999} for more information.	
\end{enumerate}

There should be no danger of confusion of $\mathcal A_n$
or $\mathcal S_n$ with the notation for a point stabilizer, where
also $n$ as a point might appear in the lower index, or 
with an automaton from an automaton family indexed by~$n$.
For the former, this should not arise since 
groups are not denoted by calligraphic letters in the present work,
and for the latter the context should make it clear whether
we refer to an automaton or the alternating group.
	
% LATA beweis review nachgucken: als fussnote adden
\begin{proposition}
	\label{sync_max:prop:prim_implies_transitive}
	A primitive permutation group $G$ 
	on a domain containing at least three points
	is transitive.
\end{proposition}
\begin{proof}
	Suppose~$G$ is not transitive.
	Since $X$ contains at least three points,
	there exist (not necessarily distinct) $a,b \in X$
	such that the set $\Delta = \Orbit{a}{G} \cup \Orbit{b}{G}$
	is a proper subset of~$X$ and contains at least two points.
	Now, $\Delta$ is a block, and so~$G$
	is not primitive.\footnote{Alternatively, for the trivial group every subset of $[n]$ is a block.
		So, for $n \ge 3$, a primitive $G$ must be non-trivial.
		Hence, there exist $a \in [n]$
		and $g \in G$ with $\applyfunc{a}{g} \ne a$.
		Then $|\Orbit{a}{G}| > 1$
		and as $\Orbit{a}{G}$ is a block 
		we find $\Orbit{a}{G} = [n]$, i.e., $G$ is transitive.}
\end{proof}

As defined here, every permutation group on $[2]$ is primitive, even 
the intransitive trivial permutation group. To exclude this special case,
some authors prefer to include transitivity into the definition of primitivity.
By \cref{sync_max:prop:prim_implies_transitive} this makes no difference
for permutation groups on domains containing at least three points.

\begin{example}\label{sync_max:ex:primitivegrps}
	\begin{enumerate}
		\item By \cref{prelim:ex:trans-block},
		for prime $n$ and transitive $G \le_G \mathcal S_n$, each block is trivial
		and so $G$ is primitive. 
		
		\item \label{sync_max:ex:primitivegrps:2} The symmetric group $\mathcal S_n$ is primitive.
		For suppose $A \subseteq [n]$ is any subset
		such that $a,b \in A$ and $c \notin A$ for distinct $a,b,c \in [n]$.
		Then, for $t = (b ~ c)$, we have $a \in \applyfunc{A}{t} \cap A$
		and $c \in t(A) \setminus A$, which shows that $A$
		is not a block. 
	\end{enumerate}
\end{example} 

If $\sim$ is a $G$-congruence, then we have the following equivalence for all $g \in G$ and $a,b \in X$:
\begin{equation}\label{eqn:G-congruence}
	a \sim b \Leftrightarrow \applyfunc{a}{g} \sim \applyfunc{b}{g}.
\end{equation}

Here, we are exclusively concerned with finite permutation domains.
So, from now on, in the remaining examples and results of the present work we always take $X = [n]$.

\begin{example}
	\label{prelim:permgrp:ex:Kleinfourgroup} 
	The \emph{Klein four-group} $G = \{ (), (0~1)(2~3), (0~2)(1~3), (0~3)(1~2) \}$
	is a regular subgroup of $\mathcal S_4$	
	where $\{0,1\}$ and $\{2,3\}$ form a system of blocks. \IN{Klein four-group}
\end{example} 

\begin{example}\label{prelim:ex:trans-block}
	Let $G \le_G \mathcal S_n$ be transitive and $\Delta \subseteq [n]$ be some block. 
	Then the sets $\applyfunc{\Delta}{g}$ for $g \in G$
	are equal-sized blocks partitioning $[n]$. So,~$n$ is a multiple of~$|\Delta|$.
\end{example}

First, we characterize the $G$-congruences
generated by 
a permutation cyclically permuting its domains.
% a cycle moving all points.
A related result appears in~\cite[Lem. 2]{HeathIKS09}.

\begin{lemma}
	\label{lem:congruence_on_cyclic_group}
	Set $g = (0 ~ 1 ~ 2 ~ \cdots ~ n - 1)_{[n]}$ and $G = \langle \{ g \} \rangle_G$. Then for every $G$-congruence~$\sim$
	there exists a divisor $r$ of $n$ such that,  for all $a,b \in [n]$,
	\begin{equation}\label{sync_max:eqn:circular:cyclic_congruence}
		a \sim b \Leftrightarrow a \equiv b \pmod{r}.
	\end{equation}
	Conversely, each divisor $r$ of $n$ defines a $G$-congruence on $[n]$
	by \cref{sync_max:eqn:circular:cyclic_congruence}.
\end{lemma}
\begin{proof} % oder i < j mit r = j - i minimal wählen und dann zeigen.
	First, suppose $\sim$ is a $G$-congruence.
	Let $0 < r \le n$ be the smallest number such that $g^r(0) \sim 0$.
	Inductively, we find $0 \sim g^{r \cdot i}(0)$ for all $i \ge 0$.
	Since $g^{i}(0) \equiv i \pmod{n}$ for $i \ge 0$,
	it is $\{ x \in [n] \mid x \equiv 0 \pmod{r} \} \subseteq [0]_{\sim}$.
	Assume $0 \sim g^{r \cdot i + d}(0)$ for some $i \ge 0$ and $0 \le d < r$.
	Since $0 \sim g^{r \cdot i}(0)$ we find $g^d(0) \sim g^{r\cdot i + d}(0)$.
	So, $g^{d}(0) \sim 0$, which implies $d = 0$ by
	the minimal choice of~$r$.
	Hence, $[0]_{\sim} \subseteq \{ x \in [n] \mid x \equiv 0 \pmod{r} \}$.
	Now, for $0 \le a < b < n$ we get
	\begin{align*}
		a \equiv b \pmod{r} 
		& \Leftrightarrow 0 \equiv b - a \pmod{r} \\
		& \Leftrightarrow 0 \sim b - a 
		\Leftrightarrow g^a(0) \sim g^a(b - a) 
		\Leftrightarrow a \sim b,		
	\end{align*}
    which shows \cref{sync_max:eqn:circular:cyclic_congruence}.
	Conversely, as $g(x) = ( x + 1 ) \bmod{n}$ for $x \in [n]$
    the relation $a \equiv b \pmod{r}$ yields a $G$-congruence when $n$ is a multiple of $r$. 	
\end{proof}

Next, we collect some properties of semiregular permutation groups.

\begin{proposition} \label{prelim:groups:prop:semiregular-size}
	Let $G \le \mathcal S_n$ be semiregular.
	Then $|G| \le n$, and we have $|G| = n$ if and only if $G$ is regular.
\end{proposition}
\begin{proof}
	Pick any point $a \in [n]$.
	Consider the map $\varphi \colon G \to [n]$
	given by $\varphi(g) = g(a)$.
	Since~$G$ is semiregular, this map is injective.
	Thus $|G| \le n$.
	Furthermore, $\varphi$ is surjective
	if and only if $G$ is transitive,
	and~$\varphi$ is surjective precisely if $|G| = n$.
	This yields that $|G| = n$ if and only if~$G$ is
	regular.
\end{proof}

%Generalizing \crefitem{prelim:permgrp:ex:singleperm}{prelim:permgrp:ex:singleperm:2}, 
Next, we characterize semiregularity by cycle lengths.
%next we state a relation between semiregularity and cycle lengths.

\begin{proposition}
 Let $G \le_G \mathcal S_n$. Then $G$
 is semiregular if and only if 
 for each $g \in G$
 all cycles contained in $g$ have the same length. %and $g$ has no fixpoints.
\end{proposition} 
\begin{proof}
	First, suppose $G$ is not semiregular.
	Then there exist $a \in [n]$ and  $g \in \Stab{a}{G} \setminus \{ \operatorname{id}_{[n]}\}$.
	As $a \in \fix{g}$, $g$ contains a cycle of length one
	and, as $g \ne \operatorname{id}_{[n]}$,
	it contains a cycle of length greater than one.	 
	Conversely, suppose there exists $g \in G$ containing two cycles
	of different lengths.
	Then 
	$g^i(a) = a$ and $g^{j}(b) = b$ with $g(a) \ne a$ and $g(b) \ne b$
	for minimal $0 < i < j < n$
	and some $a,b \in [n]$, giving $g^i \in \Stab{a}{G}$ but $g^i \ne \operatorname{id}_{[n]}$
	as $g^i(b) \ne b$. So, $G$ is not semiregular.
\end{proof}

For a semiregular $G \le_G \mathcal S_X$ and $a,b \in X$, 
since  $g(a) = g'(a)$ implies $g(b) = g'(b)$
for $g \in G$,
we can define a bijection 
$h \colon \Orbit{a}{G} \to \Orbit{b}{G}$
by setting $h(c) = g(b)$ for $g \in G$ with $\applyfunc{a}{g} = c$, i.e.,
the edges of the orbital 
digraph $\{(\applyfunc{a}{g}, \applyfunc{b}{g}) \mid g \in G\}$  
describe a bijection.
If~$G$ is transitive, $h$ is a permutation on~$X$.
In this case, since, by \cref{prelim:conj},
a permutation mapping $a$ to $b$
maps $g(a)$ to $g(b)$ if and only if $h^g = h$,
the permutation described by the tuples $\{(\applyfunc{a}{g}, \applyfunc{b}{g}) \mid g \in G\}$ 
is characterized as the permutation~$h$
with $h(a) = b$ and $hg = gh$ for all $g \in G$.
Thus, if~$G$ is regular and commutative,
it contains a permutation described by
the tuples $\{(\applyfunc{a}{g}, \applyfunc{b}{g}) \mid g \in G\}$.

\begin{example} \label{prelim:permgrp:ex:singleperm}
	Let $G = \langle \{ t\} \rangle_G$ for some permutation 
	$t \colon [n] \to [n]$.	Then,
	
	\begin{enumerate}
		\item The orbits of~$G$ are the sets $\{ t^i(x) \mid i \ge 0\}$,
		and the non-singleton sets of this form
		equal the points moved by some non-identity
		cycle induced by~$t$.
			
		\item $t$ is an cycle if and only if $G$
		has at most one non-singleton orbit,
		
		\item \label{prelim:permgrp:ex:singleperm:2} $G$ is semiregular if and only if every cycle in $t$ 
		has the same length,
		
		%\item The following are equivalent: $G$ is regular, $G$ is transitive,
		%and $t$ cyclically permutes $[n]$.
		
		\item $G$ is regular $\Leftrightarrow$ $G$ is transitive $\Leftrightarrow$ $t$ cyclically permutes~$[n]$.
		
		%\item $G$ is regular if and only if $t$ cyclically permutes $[n]$.
	\end{enumerate} 
\end{example} 

\section{Transitive Permutation Groups}
\label{prelim:permgrp:sec:transitive}

The following three subsections collect results on permutation groups related to transitivity.
\Cref{prelim:perm:subsec:sep-sets-idempotents} treats separation of subsets by a transitive
group and shows which idempotents arise in the transformation semigroup generated by a
transitive group together with a one-point singular.
\Cref{prelim:perm:subsec:blocks-stabs} establishes a correspondence between the subgroups
containing the stabilizer of a point~$a$ and the blocks containing~$a$, under which cosets
correspond to images of blocks.
\Cref{prelim:perm:subsec:mapping-orbit} relates conditions on how a group maps $2$-sets to one
another to the size of orbits, to transitivity, and to $2$-homogeneity.

\subsection{Separating Sets and Idempotents}
\label{prelim:perm:subsec:sep-sets-idempotents}

In transitive permutation groups we can separate subsets in the following way
(see~\cite[p.~130]{ArCamSt2017} for related discussion).

\begin{theorem}[\cite{ArCamSt2017,birch2009,neumann75b}]
	\label{thm:separating_neumann}
	Let $G \le_G \mathcal S_n$ be transitive and let $A,B \subseteq [n]$ with $|A| \cdot |B| < n$.
	Then there exists $g \in G$ such that $g(A) \cap B = \emptyset$.
\end{theorem}

Our next results let us construct certain idempotent one-point singulars.

\begin{proposition} \label{prelim:prop:one-point-singular-power}
	Let $t \colon [n] \to [n]$ be one-point singular with $\dupl{t} = \{c\}$.
	Then there exists an idempotent $h \in \langle \{t\} \rangle_S$
	and we have the following:
	\begin{enumerate}
		\item \label{prelim:prop:one-point-singular-power:1} If $t(c) = c$, then~$h$ is one-point
		singular with $\dupl{h} = \{c\}$ and $t^{-1}(c) \subseteq h^{-1}(c)$.
		
		\item \label{prelim:prop:one-point-singular-power:2} If~$t$ has defect one and the point
		excluded by~$t$ lies in $t^{-1}(c)$, say $\excl{t} = \{a\}$ and $t^{-1}(c) = \{a,b\}$,
		then~$h$ is the simple idempotent mapping~$a$ to~$b$.
	\end{enumerate}
\end{proposition}
\begin{proof}
	By \cref{prelim:lem:finite-semigrp-idempotent-power} such an~$h$ exists; write $h = t^i$
	with $i > 0$.
	
	\ref{prelim:prop:one-point-singular-power:1}: We show $\dupl{t^j} = \{c\}$ for $j \ge 1$ by
	induction, the case $j = 1$ being clear.
	For $j \ge 2$ write $t^j = \concat{t^{j-1}}{t}$, so that \cref{prelim:eqn:dupl0} yields
	$\dupl{t^j} = \dupl{t_{|t^{j-1}([n])}} \cup t(\dupl{t^{j-1}})$, where the first set is
	contained in $\dupl{t} = \{c\}$ and the second equals $t(\{c\}) = \{c\}$.
	Hence $\dupl{h} = \{c\}$.
	Moreover $t(x) = c$ gives $h(x) = t^{i-1}(c) = c$, so $t^{-1}(c) \subseteq h^{-1}(c)$.
	
	\ref{prelim:prop:one-point-singular-power:2}: We have $t([n]) = [n] \setminus \{a\}$, and
	$t(a) = t(b) = c$ with $b \ne a$ gives $t([n] \setminus \{a\}) = t([n]) = [n] \setminus \{a\}$.
	So~$t$ maps the finite set $[n] \setminus \{a\}$ onto itself and thus permutes it.
	Hence~$h$ restricts to an idempotent permutation of $[n] \setminus \{a\}$, that is, to the
	identity, and $h(a) = t^{i-1}(t(a)) = t^{i-1}(t(b)) = h(b) = b$.
\end{proof}

\begin{remark}
	The hypotheses of the two items differ in which point is required to lie in $t^{-1}(c)$:
	in \cref{prelim:prop:one-point-singular-power:1} it is the duplicated point~$c$ itself, and
	in \cref{prelim:prop:one-point-singular-power:2} it is the excluded point~$a$.
	This accounts for the stronger conclusion in the second case.
	There, $t$ maps $[n] \setminus \{a\}$ onto itself, so no point outside $t^{-1}(c)$ is ever
	collapsed by a power of~$t$ and we obtain $h^{-1}(c') = t^{-1}(c)$ for the point $c'$ with
	$\dupl{h} = \{c'\}$.
	Under \cref{prelim:prop:one-point-singular-power:1} alone, the inclusion
	$t^{-1}(c) \subseteq h^{-1}(c)$ can be strict, as passing to powers may collapse further
	points: for $n = 4$ and~$t$ given by $t(0) = t(1) = t(2) = 0$ and $t(3) = 2$ we have
	$c = 0$ and $t^{-1}(0) = \{0,1,2\}$, whereas $t^2$ is the constant map at~$0$, so that
	$h^{-1}(0) = [4]$.
\end{remark}

\begin{corollary} \label{prelim:lem:transitive-grp-one-point-singular}
	Let $G \le_G \mathcal S_n$ be a transitive permutation group and let $t \colon [n] \to [n]$
	be one-point singular.
	Then for each $a \in t^{-1}(\dupl{t})$ there exists an idempotent one-point singular
	$h \in \langle G \cup \{t\} \rangle_S$ with $\dupl{h} = \{a\}$ and
	$t^{-1}(\dupl{t}) \subseteq h^{-1}(a)$.
\end{corollary}
\begin{proof}
	Write $\dupl{t} = \{c\}$ and pick $g \in G$ with $g(c) = a$, which is possible as~$G$ is transitive.
	As~$g$ is a bijection, $\concat{t}{g}$ is one-point singular with
	$\dupl{\concat{t}{g}} = \{a\}$ and $(\concat{t}{g})^{-1}(a) = t^{-1}(\dupl{t})$.
	The latter contains~$a$, so $\concat{t}{g}$ fixes~$a$ and
	\cref{prelim:prop:one-point-singular-power:1} of
	\cref{prelim:prop:one-point-singular-power}, applied to $\concat{t}{g}$, gives an idempotent
	$h \in \langle \{\concat{t}{g}\} \rangle_S \subseteq \langle G \cup \{t\} \rangle_S$ as
	claimed.
\end{proof}

\begin{corollary}
	\label{sync_words:lem:idempotent}
	Let $G \le_G \mathcal S_n$ be a transitive permutation group and let $t \colon [n] \to [n]$
	be of defect one with $t(a) = t(b)$ for distinct $a,b \in [n]$.
	Then there exists a simple idempotent $h \in \langle G \cup \{t\} \rangle_S$
	with $h(a) = h(b) = b$.
\end{corollary}
\begin{proof}
	Write $\excl{t} = \{e\}$ and pick $g \in G$ with $g(e) = a$, which is possible as~$G$ is transitive.
	As~$g$ is a bijection, $\concat{t}{g}$ has defect one with
	$\excl{\concat{t}{g}} = \{a\}$ and $(\concat{t}{g})^{-1}(\dupl{\concat{t}{g}}) =
	t^{-1}(\dupl{t}) = \{a,b\}$.
	Now \cref{prelim:prop:one-point-singular-power:2} of
	\cref{prelim:prop:one-point-singular-power}, applied to $\concat{t}{g}$, gives an idempotent
	$h \in \langle \{\concat{t}{g}\} \rangle_S \subseteq \langle G \cup \{t\} \rangle_S$
	mapping~$a$ to~$b$.
\end{proof}

Applying \cref{sync_words:lem:idempotent} to the pair $(a,b)$ and to the pair $(b,a)$ shows that
for $t$ of defect one with $t(a) = t(b)$ the semigroup $\langle G \cup \{t\} \rangle_S$
contains \emph{both} simple idempotents identifying $a$ and $b$.

\begin{example}
	Let $g = (0 ~ 1 ~ 2 ~ 3)$ and let $t \colon [4] \to [4]$ be given by
	$t(0) = 1$, $t(2) = 3$ and $t(1) = t(3) = 2$.
	Then $gt$ permutes $\{0,2,3\}$ and $g^3 t$ permutes $\{0,1,2\}$.
	Here $g^3 t$ is already the simple idempotent mapping $3$ to $1$,
	and $(gt)^2$ is the simple idempotent mapping $1$ to $3$.
\end{example}

\subsection{Blocks and Point Stabilizers}
\label{prelim:perm:subsec:blocks-stabs}

We next establish a correspondence between the subgroups containing the stabilizer of a point~$a$
and the blocks containing~$a$ in \cref{prelim:thm:blocks-stabs:2} of \cref{prelim:thm:blocks-stabs}. This also appears
in~\cite[Thm.~1.5A]{dixon1996permutation}.

\begin{theorem}\label{prelim:thm:blocks-stabs}
	Let $G \le_G \mathcal S_n$ be transitive and $a \in [n]$. Then the following hold.
	\begin{enumerate}
		\item \label{prelim:thm:blocks-stabs:2} The maps
		\[
		H \mapsto \Orbit{a}{H}
		\qquad\text{and}\qquad
		\Delta \mapsto \{ g \in G \mid g(\Delta) = \Delta \}
		\]
		are mutually inverse bijections between the subgroups $H$ with $\Stab{a}{G} \le_G H \le_G G$
		and the blocks $\Delta$ of $G$ with $a \in \Delta$.
		
		\item \label{prelim:thm:blocks-stabs:3} For $\Stab{a}{G} \le_G H \le_G G$,
		$\Delta = \Orbit{a}{H}$ and $g_1,g_2 \in G$ we have
		\[
		g_1 H = g_2 H \iff g_1(\Delta) = g_2(\Delta),
		\]
		and hence $gH = \{ g' \in G \mid g'(\Delta) = g(\Delta) \}$.
	\end{enumerate}
\end{theorem}
\begin{proof}
	\ref{prelim:thm:blocks-stabs:2}:
	Let $\Stab{a}{G} \le_G H \le_G G$ and $\Delta = \Orbit{a}{H}$.
	We show
	\begin{equation}\label{prelim:eqn:block-stab-key}
		\{ g \in G \mid g(\Delta) \cap \Delta \ne \emptyset \} = H .
	\end{equation}
	If $g \in H$ then $g(\Delta) = \Delta$, since $\Delta$ is an $H$-orbit; in particular the
	left-hand side contains $H$.
	Conversely, suppose $g(h_2(a)) = h_1(a)$ for some $h_1,h_2 \in H$.
	Then $h_1^{-1} g h_2 \in \Stab{a}{G} \le_G H$ and therefore $g \in h_1 H h_2^{-1} = H$.
	This proves~\eqref{prelim:eqn:block-stab-key}, and, since $g \in H$ implies $g(\Delta) = \Delta$,
	it shows both that $\Delta$ is a block and that the second map returns $H$ when applied to $\Delta$.
	
	Now let $\Delta$ be a block with $a \in \Delta$ and set $H = \{ g \in G \mid g(\Delta) = \Delta \}$,
	which is a subgroup.
	If $g \in \Stab{a}{G}$ then $a \in g(\Delta) \cap \Delta$, so $g(\Delta) = \Delta$ and $g \in H$;
	hence $\Stab{a}{G} \le_G H$.
	The inclusion $\Orbit{a}{H} \subseteq \Delta$ is immediate from the definition of $H$.
	Conversely let $b \in \Delta$. By transitivity $b = g(a)$ for some $g \in G$, so
	$b \in g(\Delta) \cap \Delta$ and thus $g(\Delta) = \Delta$, i.e.\ $g \in H$ and $b \in \Orbit{a}{H}$.
	Hence $\Delta = \Orbit{a}{H}$, so the first map returns $\Delta$, and the two maps are
	mutually inverse.
	
	\ref{prelim:thm:blocks-stabs:3}:
	By~\eqref{prelim:eqn:block-stab-key} we have $g \in H$ if and only if $g(\Delta) = \Delta$. Thus
	\[
	g_1 H = g_2 H
	\iff g_2^{-1} g_1 \in H
	\iff g_2^{-1} g_1 (\Delta) = \Delta
	\iff g_1(\Delta) = g_2(\Delta).
	\]
	That $gH = \{ g' \in G \mid g'(\Delta) = g(\Delta) \}$
	follows since $g' \in gH$ precisely if $g'H = gH$.
\end{proof}

\begin{remark}
	By \cref{prelim:thm:blocks-stabs:3} of \cref{prelim:thm:blocks-stabs} the map $\varphi(gH) = g(\Delta)$
	is a well-defined bijection from the cosets of $H$ onto the block system $\{ g(\Delta) \mid g \in G\}$
	generated by
	$\Delta = \Orbit{a}{H}$, and $\varphi(h \cdot gH) = h(\varphi(gH))$ for $h \in G$.
	Thus subgroups containing $\Stab{a}{G}$ together with their cosets correspond to blocks
	containing $a$ together with their block systems, and,
	for $g \in G$ and $h \in G$, we have
	\begin{equation}\label{eqn:congruences-subgroups-blocks}
		h \in gH \iff h(a) \in g(\Delta).
	\end{equation}
	As an example, let $G$ be as in \cref{lem:congruence_on_cyclic_group}.
	Then $G$ can be identified with the additive group of residues modulo $n$ via
	$g^{i} \mapsto [i]_{\equiv_n}$.
	Since every subgroup of that group is the group of multiples of some divisor $r$ of $n$
	\cite[Sec.~1.4]{Kurzweil04}, \cref{eqn:congruences-subgroups-blocks} 
	also implies \cref{lem:congruence_on_cyclic_group}.
\end{remark}

\begin{remark}
	\Cref{prelim:thm:blocks-stabs} is an instance of the standard correspondence between
	transitive $G$-actions and actions on the cosets of a subgroup: a system of blocks induces a
	$G$-action for which~$H$ is a point stabilizer, and the theorem states that the actions
	arising in this way are precisely those on the cosets of a subgroup containing a point
	stabilizer.
	We refer to~\cite{cameron_1999} for the formalisation via $G$-actions, which we do not
	introduce here, and note that by our conventions permutations act on the left, whereas
	in~\cite{cameron_1999} they act on the right; accordingly, our statements concern left cosets
	and multiplication from the left.
\end{remark}

\IN{orbit-stabilizer property}
\begin{remark}\label{prelim:rem:orbit-stab}
	Specialising \crefitem{prelim:thm:blocks-stabs}{prelim:thm:blocks-stabs:3} to $H = \Stab{a}{G}$,
	the left cosets of $\Stab{a}{G}$ are in bijection with the points of $[n]$.
	As the cosets partition $G$ and all have size $|\Stab{a}{G}|$, we obtain
	$|G| = n \cdot |\Stab{a}{G}|$, the \emph{orbit-stabilizer property}
	\cite{cameron_1999,dixon1996permutation,Kurzweil04,wielandt1964finite}.
	This implies that a transitive $G \le_G \mathcal S_n$ is regular if and only if $|G| = n$.
\end{remark}

\begin{corollary}
	\label{sync_max:cor:point_stabs}
	Let $G \le_G \mathcal S_n$ be transitive, $a \in [n]$ and $\Stab{a}{G} \le_G H \le_G G$.
	Set $\Delta = \Orbit{a}{H}$. Then:
	\begin{enumerate}
		\item \label{sync_max:cor:point_stabs:1} For $g \in G$ we have $g \notin H$ if and only if $a \notin g(\Delta)$.
		\item \label{sync_max:cor:point_stabs:Htrans} If $H \ne G$ then $H$ does not act transitively.
	\end{enumerate}
\end{corollary}
\begin{proof}
	\ref{sync_max:cor:point_stabs:1}: 
	Since $\Delta$ is a block by \cref{prelim:thm:blocks-stabs}, $a \in g(\Delta)$
	is equivalent to $g(\Delta) = \Delta$ for $g \in G$.
	Since $H = \{ g\in G \mid g(\Delta) = \Delta \}$
	by \cref{prelim:thm:blocks-stabs:2} of \cref{prelim:thm:blocks-stabs},
	$g \in G \setminus H$ if and only if $a \notin g(\Delta)$.
	
	\ref{sync_max:cor:point_stabs:Htrans}: By
	\crefitem{prelim:thm:blocks-stabs}{prelim:thm:blocks-stabs:2} the assignment
	$H \mapsto \Orbit{a}{H}$ is injective, and $\Orbit{a}{G} = [n]$ by transitivity of $G$.
	Hence $\Orbit{a}{H} = [n]$ forces $H = G$.
\end{proof}

\subsection{Mapping $2$-Sets, Orbit Sizes, and Transitivity}
\label{prelim:perm:subsec:mapping-orbit}

We now relate mapping properties of $2$-sets to orbit sizes and transitivity.
Two technical statements do most of the work.
\Cref{prelim:grp:prop:map-sets-outside-set-A} treats images of the form $\{a,y\} \mapsto \{b,y\}$,
where a point of the source is retained, and
\cref{prelim:prop:single-two-set} settles, for a single target $2$-set, when it is the image of a
$2$-set containing a fixed point~$a$ under a map subject to a side condition.
Quantifying the latter over all targets yields \cref{prelim:cor:two-set-global}, which relates
several such mapping properties to the size of $\Orbit{a}{G}$, to transitivity, and to the fixed
points of the point stabiliser of~$a$.
From these we obtain characterisations of transitivity in \cref{prelim:thm:trans} and
sufficient conditions in \cref{prelim:prop:suff-transitive}.
Finally, \cref{prelim:prop:2homogeneous-pair} determines which mapping properties force
$2$-homogeneity, i.e., that any two $2$-sets can be mapped to one another.
Before that we recall a result of Livingstone and Wagner~\cite{Livingstone1965} relating
$k$-homogeneity to $k$-transitivity.

\begin{theorem}[Livingstone and Wagner~\cite{Livingstone1965}]
	\label{prelim:thm:livingstone}
	Let $G \le_G \mathcal S_n$ be $k$-homogeneous with $k \ge 2$ and $2k \le n$.
	Then $G$ is $(k-1)$-transitive, and for $k \ge 5$ even $k$-transitive.
\end{theorem}

The argument from the following lemma 
is taken from~\cite[p.~533]{mcalister1998semigroups}, where it is used to deduce
that a $2$-homogeneous group is transitive.

\begin{lemma} \label{prelim:grp:lem:three-points-mapping}
	Let $G \le_G \mathcal S_n$ and $a,b,c \in [n]$.
	If $g(\{a,b\}) = \{c,b\}$ for some $g \in G$,
	then $c \in \Orbit{a}{G}$.	
\end{lemma}
\begin{proof}
    We have $g(a) = c$ or $g(a) = b$.
    In the former case, $c \in \Orbit{a}{G}$.
    In the latter case, we have  $g(b) = c$ and thus $g^2(a) = c$.
    So also $c \in \Orbit{a}{G}$. 
\end{proof}

\begin{proposition} \label{prelim:grp:prop:map-sets-outside-set-A}
	Let $G \le_G \mathcal S_n$, let $a \in [n]$ and let $A \subseteq [n]$.
	Suppose that for every $b \in A$ there are $y \in [n]$ and $g \in G$ with
	$g(\{a,y\}) = \{b,y\}$.
	Then
	\[
	A \cup \{a\} \subseteq \Orbit{a}{G}.
	\]
\end{proposition}
\begin{proof}
	By \cref{prelim:grp:lem:three-points-mapping}
	we have $b \in \Orbit{a}{G}$ for every $b \in A$,
	and $a \in \Orbit{a}{G}$ holds anyway.
\end{proof}

\begin{remark}
	The method of proof of \cref{prelim:grp:lem:three-points-mapping}
	actually established $c \in B$ for
	\[
	B = \{g^2(a) \mid g \in G\} \cup
	\{g(a) \mid g \in G \text{ with } g(b) = b \text{ for some } b \ne a \}.
	\]
	Thus $A$ from \cref{prelim:grp:prop:map-sets-outside-set-A} is also contained
	in this set.
	For $A = [n]$ this yields that, for fixed $a \in [n]$, for each $c \in [n]$ there exist
	$b \in [n]$ and $g \in G$ such that $g(\{a,b\}) = \{c,b\}$ if and only if $[n] = B$.
	As $B \subseteq \Orbit{a}{G}$, this forces $G$ to be transitive.
	For example, we have $[n] = B$ for transitive groups of odd order.
	Indeed, on a finite group the squaring map $g \mapsto g^2$ is surjective if and only if
	it is injective, and it fails to be injective precisely when there is an involution, i.e.,
	a non-identity element whose square equals the identity.
	By Cauchy's theorem an involution exists whenever $|G|$ is even, and conversely an
	involution has order~$2$, so that $2$ divides $|G|$.
	Hence for odd $|G|$ we get $\{g^2 \mid g \in G\} = G$ and therefore
	$\{g^2(a) \mid g \in G\} = \Orbit{a}{G} = [n]$.
	Note also that for transitive $G$ the condition does not depend on the choice of $a$:
	conjugation by an element $h \in G$ with $h(a) = a'$ carries the two sets
	for $a$ onto those for $a'$.
\end{remark}

\begin{remark}
	We do not need this fact, but it is worth noting
	that, for fixed $a,b \in [n]$,
	if, for every $c \in [n]$,
	we have  $g(\{a,b\}) = \{a,c\}$ or $g(\{a,b\}) = \{b,c\}$ 
	for some $g \in G$, 
	then $[n] = \Orbit{a}{G} \cup \Orbit{b}{G}$.
	This follows since then  $c \in \Orbit{a}{G} \cup \Orbit{b}{G}$ for 
	each $c \in [n]$.
\end{remark}

\begin{proposition} \label{prelim:prop:single-two-set}
	Let $G \le_G \mathcal S_n$, let $a \in [n]$ and let $\{c,d\} \subseteq [n]$ be a $2$-set.
	\begin{enumerate}
		\item \label{prelim:prop:single-two-set:T} For $T \subseteq [n]$ there are $g \in G$ and
		$b \in [n]$ with $g(\{a,b\}) = \{c,d\}$ and $g(a) \in T$ if and only if
		$\{c,d\} \cap \Orbit{a}{G} \cap T \ne \emptyset$.
		
		\item \label{prelim:prop:single-two-set:1} There are $g \in G$ and $b \in [n]$ with
		$g(\{a,b\}) = \{c,d\}$ and $g(b) \ne b$ if and only if there is $g \in G$ with $g(a) = c$
		and $g(d) \ne d$, or with $g(a) = d$ and $g(c) \ne c$.
	\end{enumerate}
\end{proposition}
\begin{proof}
	\ref{prelim:prop:single-two-set:T}: Given such $g$ and $b$, we have
	$g(a) \in g(\{a,b\}) = \{c,d\}$ as well as $g(a) \in \Orbit{a}{G}$ and $g(a) \in T$.
	Conversely, for $x \in \{c,d\} \cap \Orbit{a}{G} \cap T$ and $y$ the other element of
	$\{c,d\}$, choose $g \in G$ with $g(a) = x$ and put $b = g^{-1}(y)$;
	then $g(a) \in T$ and $g(\{a,b\}) = \{x,y\} = \{c,d\}$.
	
	\ref{prelim:prop:single-two-set:1}: Given such $g$ and $b$, we have $g(a) \in \{c,d\}$;
	say $g(a) = c$, so that $g(b) = d$ and $b = g^{-1}(d)$, and $g(b) \ne b$ gives
	$d \ne g^{-1}(d)$, i.e., $g(d) \ne d$.
	Conversely, for $g(a) = c$ and $g(d) \ne d$ put $b = g^{-1}(d)$;
	then $g(\{a,b\}) = \{c,d\}$ and $g(b) = d \ne b$, as $b = d$ would give $g(d) = d$.
	Interchanging $c$ and $d$ gives the remaining cases.
\end{proof}

\begin{corollary} \label{prelim:cor:two-set-global}
	Let $G \le_G \mathcal S_n$ with $n \ge 2$ and $a \in [n]$.
	\begin{enumerate}
		\item \label{prelim:cor:two-set-global:1} \label{prelim:lem:2sets-transitive}
		For every $2$-set $\{c,d\} \subseteq [n]$ there are $g \in G$ and $b \in [n]$ with
		$g(\{a,b\}) = \{c,d\}$ if and only if $|\Orbit{a}{G}| \ge n - 1$.
		
		\item \label{prelim:cor:two-set-global:2} For every $2$-set $\{c,d\} \subseteq [n]$ there are
		$g \in G$ and $b \in [n]$ with $g(\{a,b\}) = \{c,d\}$ and $g(a) \ne a$ if and only if~$G$ is
		transitive.
		
		\item \label{prelim:cor:two-set-global:3} For every $b \in [n] \setminus \{a\}$ there are
		$g \in G$ and $c \in [n]$ with $g(\{a,c\}) = \{a,b\}$ and $g(a) \ne a$ if and only if~$G$ is
		transitive.
		
		\item \label{prelim:cor:two-set-global:4} For every $b \in [n] \setminus \{a\}$ 
		there are
		$g \in G$ and $c \in [n]$ 
		with $g(\{a,c\}) = \{a,b\}$ and $g(c) \ne c$ 
		if and only if every point fixed by all 
		of $\Stab{a}{G}$ lies in $\Orbit{a}{G}$.
	\end{enumerate}
\end{corollary}
\begin{proof}
	\ref{prelim:cor:two-set-global:1}: By
	\cref{prelim:prop:single-two-set:T} of \cref{prelim:prop:single-two-set} with $T = [n]$ the condition
	states that every $2$-set meets $\Orbit{a}{G}$, which holds if and only if
	$|[n] \setminus \Orbit{a}{G}| \le 1$, as a $2$-set contained in that complement exists
	precisely when it has at least two elements.	
	
	\ref{prelim:cor:two-set-global:2} and \ref{prelim:cor:two-set-global:3}: By
	\cref{prelim:prop:single-two-set:T} of \cref{prelim:prop:single-two-set} with
	$T = [n] \setminus \{a\}$, both conditions state that the $2$-sets in question meet
	$\Orbit{a}{G} \cap T = \Orbit{a}{G} \setminus \{a\}$.
	
	For a target $\{a,b\}$ with $b \ne a$ this says that $b \in \Orbit{a}{G}$, so
	\cref{prelim:cor:two-set-global:3} implies that $G$ is transitive (For a direct argument,
	for the target set $\{a,b\}$
    we have $g(a) \in \{a,b\}$ with $g(a) \ne a$, which gives $g(a) = b$.).
	If $G$ is transitive, then $\Orbit{a}{G} \setminus \{a\} = [n] \setminus \{a\}$, which is met
	by every $2$-set, giving \cref{prelim:cor:two-set-global:2}.
	Finally, \cref{prelim:cor:two-set-global:2} implies \cref{prelim:cor:two-set-global:3}, as the
	latter only concerns the targets $\{a,b\}$.
	Hence \cref{prelim:cor:two-set-global:2}, \cref{prelim:cor:two-set-global:3} and transitivity
	of~$G$ are equivalent, which proves both stated equivalences.
		
	\ref{prelim:cor:two-set-global:4}: Consider some $b \in [n] \setminus\{a\}$. By
	\cref{prelim:prop:single-two-set:1} of \cref{prelim:prop:single-two-set} applied to $\{a,b\}$,
	such $g$ and $c$ exist if and only if there is $g \in G$ with $g(a) = a$ and $g(b) \ne b$,
	or with $g(a) = b$, i.e., if and only if some element of $\Stab{a}{G}$ moves~$b$ or
	$b \in \Orbit{a}{G}$.
	This is equivalent to the condition that if the considered~$b$
	is fixed by all elements from $\Stab{a}{G}$,
	then it must be in $\Orbit{a}{G}$.
	Demanding this for every $b \in [n] \setminus \{a\}$ says precisely that a point fixed by
	all of $\Stab{a}{G}$ lies in $\Orbit{a}{G}$.
\end{proof}

\begin{theorem} \label{prelim:thm:trans}
	Let $G \le_G \mathcal S_n$ with $n \ge 2$ and $a \in [n]$. Then the following are equivalent.
	\begin{enumerate}
		\item \label{prelim:thm:trans:1} $G$ is transitive.
		
		\item \label{prelim:thm:trans:2} For distinct $c,d \in [n]$ there is $g \in G$ with
		$g(a) = c$ and $g(d) \ne d$, or $g(a) = d$ and $g(c) \ne c$.
		
		\item \label{prelim:thm:trans:3} Every point fixed by all of $\Stab{a}{G}$ lies in
		$\Orbit{a}{G}$, and every point of $[n] \setminus \Orbit{a}{G}$ is fixed by all of~$G$.
	\end{enumerate}
\end{theorem}
\begin{proof}
	$\ref{prelim:thm:trans:1}\Rightarrow\ref{prelim:thm:trans:2}$: Choose $g_0, g_1 \in G$
	with $g_0(a) = c$ and $g_1(a) = d$.
	If $g_0(d) \ne d$ or $g_1(c) \ne c$, we are done.
	Otherwise $h = g_1^{-1}g_0$ satisfies $h(a) = g_1^{-1}(c) = c$ and
	$h(d) = g_1^{-1}(d) = a \ne d$, the latter since $a = d$ would give
	$c = g_0(a) = g_0(d) = d$.
	
	$\ref{prelim:thm:trans:2}\Rightarrow\ref{prelim:thm:trans:3}$: By
	\cref{prelim:prop:single-two-set:1} of \cref{prelim:prop:single-two-set}, condition
	\ref{prelim:thm:trans:2} gives, for every $2$-set $\{c,d\}$, some $g \in G$ and $b \in [n]$
	with $g(\{a,b\}) = \{c,d\}$ and $g(b) \ne b$.
	Omitting $g(b) \ne b$ gives $|\Orbit{a}{G}| \ge n - 1$ by
	\cref{prelim:cor:two-set-global:1} of \cref{prelim:cor:two-set-global}, so at most one point lies
	outside $\Orbit{a}{G}$, and it forms an orbit and is thus fixed by all of~$G$.
	Restricting to the targets $\{a,b\}$ gives the remaining condition by
	\cref{prelim:cor:two-set-global:4} of \cref{prelim:cor:two-set-global}.
	
	$\ref{prelim:thm:trans:3}\Rightarrow\ref{prelim:thm:trans:1}$: 
	Let $b \in [n] \setminus \Orbit{a}{G}$.
	Then $b$ is fixed by all of~$G$, hence by all of $\Stab{a}{G}$, and so lies in
	$\Orbit{a}{G}$, a contradiction.
	Thus $\Orbit{a}{G} = [n]$.
\end{proof}

Finally, we give some sufficient conditions.

\begin{proposition} \label{prelim:prop:suff-transitive}  
	\label{prelim:lem:2sets-and-botha-transitive} % das label eigentlich auf 2.te aussage
	 \label{prelim:lem:2sets-transitive2} % das label eigenltich auf 3.te aussage
	Let $G \le_G \mathcal S_n$ with $n \ge 2$ and $a \in [n]$.
	Then~$G$ is transitive if one of the following holds:
	\begin{enumerate}
		\item \label{prelim:prop:suff-transitive:1} $|\Orbit{a}{G}| \ge n - 1$ and for
		$b \in [n] \setminus \{a\}$ we have $g(\{a,c\}) = \{a,b\}$ for some
		$c \in [n] \setminus \{b\}$ and $g \in G$.
		
		\item \label{prelim:prop:suff-transitive:2} For distinct $c,d \in [n]$ we have
		$g(\{a,b\}) = \{c,d\}$ for some $g \in G$ and $b \in [n] \setminus \{c\}$.
		
		\item \label{prelim:prop:suff-transitive:3} For every $c \in [n]$ we have
		$g(\{a,b\}) = \{c,b\}$ for some $b \in [n]$ and $g \in G$.
	\end{enumerate}
\end{proposition}
\begin{proof}
	\ref{prelim:prop:suff-transitive:1}: Such a $g$ maps $c$ to $a$ or to~$b$ with
	$c \notin \{a,b\}$, so $g(c) \ne c$, and
	\cref{prelim:cor:two-set-global:4} of \cref{prelim:cor:two-set-global} gives that every point
	fixed by all of $\Stab{a}{G}$ lies in $\Orbit{a}{G}$.
	As $|\Orbit{a}{G}| \ge n - 1$ leaves at most one point outside $\Orbit{a}{G}$, which forms an
	orbit and is fixed by all of~$G$, \cref{prelim:thm:trans} applies.
	
\ref{prelim:prop:suff-transitive:2}: By
\cref{prelim:cor:two-set-global:1} of \cref{prelim:cor:two-set-global} we have
$|\Orbit{a}{G}| \ge n - 1$, and the pair $b,a$ gives $g(\{a,c\}) = \{a,b\}$
for some $c \ne b$,	which is the condition from \cref{prelim:prop:suff-transitive:1}.
	
	\ref{prelim:prop:suff-transitive:3}: This follows from \cref{prelim:grp:prop:map-sets-outside-set-A} applied with $A = [n]$.
\end{proof}

\begin{remark}
	None of the conditions of \cref{prelim:prop:suff-transitive} is necessary for transitivity:
	for $n = 4$ the group $G = \langle (0~1~2~3) \rangle_G$ is transitive, but at $a = 0$ and
	$b = 2$ the equality $g(\{0,c\}) = \{0,2\}$ forces $c = 2$, so that
	\cref{prelim:prop:suff-transitive:1} fails, and applying
	\cref{prelim:prop:suff-transitive:2} to the pair $2,0$ would give exactly such a $c \ne 2$.
	Likewise \cref{prelim:prop:suff-transitive:3} fails for odd~$c$: in $g(\{0,b\}) = \{c,b\}$ we
	have $c \ne b$, so either $g(0) = c$ and $g(b) = b$, which for a shift forces $g$ to be the
	identity and hence $c = 0$, or $b = g(0)$ and $c = g(b) = g(g(0)) = 2k$, and both values are
	even.
\end{remark}

\begin{remark}
	The ordered reading matters in \cref{prelim:prop:suff-transitive:2} of \cref{prelim:prop:suff-transitive}.
	Demanding instead a single witness with $b \notin \{c,d\}$ is strictly stronger and is
	unsatisfiable for $n = 3$: for the $2$-set $[n] \setminus \{a\}$ there is no $b$ with
	$b \ne a$ and $b \notin [n] \setminus \{a\}$.
	By contrast, $\mathcal S_3$ satisfies the hypothesis of \cref{prelim:prop:suff-transitive:2} from \cref{prelim:prop:suff-transitive}.
	For $n = 4$ the transitive group $G = \langle (0~1~2~3) \rangle_G$ violates it at $a = 0$ and
	the $2$-set $\{0,2\}$: a witness would need $b \in \{1,3\}$, whereas $g(\{0,b\}) = \{0,2\}$
	forces $b = 2$.
\end{remark}

% Diese alternative Aussage mal noch behalten
%\begin{proposition} \label{prelim:prop:2homogeneous-pair}
%	Let $G \le_G \mathcal S_n$ with $n \ge 2$. Then the following are equivalent.
%	\begin{enumerate}
%		\item \label{prelim:prop:2homogeneous-pair:1} $G$ is $2$-homogeneous.
%		
%		\item \label{prelim:prop:2homogeneous-pair:2} There exist distinct $a,b \in [n]$ such
%		that for all $c,d \in [n]$ with $\{c,d\} \ne \{a,b\}$ there is $g \in G$ with
%		$g(\{a,c\}) = \{a,d\}$ or $g(\{b,c\}) = \{b,d\}$.
%		
%		\item \label{prelim:prop:2homogeneous-pair:3} There exist distinct $a,b \in [n]$ such
%		that every $2$-set meeting $\{a,b\}$ is of the form $g(\{a,b\})$ for some $g \in G$.
%	\end{enumerate}
%\end{proposition}
%\begin{proof}
%	%	In \ref{prelim:prop:2homogeneous-pair:2} the points $c$ and $d$ range over all of $[n]$, so
%	%	degenerate instances are admitted; as $g$ is injective, $g(\{x,x\}) = \{x\}$ is a singleton
%	%	and thus never a $2$-set.
%	\ref{prelim:prop:2homogeneous-pair:1}$\Rightarrow$\ref{prelim:prop:2homogeneous-pair:2}:
%	Let $a,b$ be any two distinct points and consider $\{c,d\} \ne \{a,b\}$.
%	Without loss of generality, suppose $a \notin \{c,d\}$.
%	Then $\{a,c\}$ and $\{a,d\}$ are $2$-sets and $2$-homogeneity gives
%	$g(\{a,c\}) = \{a,d\}$ for some $g \in G$.
%	
%	
%	\ref{prelim:prop:2homogeneous-pair:2}$\Rightarrow$\ref{prelim:prop:2homogeneous-pair:3}:
%	A $2$-set meeting $\{a,b\}$ is $\{a,b\}$ itself, the image under the identity, or is of the
%	form $\{a,e\}$ or $\{b,e\}$ with $e \notin \{a,b\}$.
%	Applying \cref{prelim:prop:2homogeneous-pair:2} to $b,e$ and to $a,e$, which is admissible as
%	$\{b,e\} \ne \{a,b\} \ne \{a,e\}$, gives $g(\{a,b\}) = \{a,e\}$ or $g(\{b,b\}) = \{b,e\}$,
%	respectively $g(\{a,a\}) = \{a,e\}$ or $g(\{b,a\}) = \{b,e\}$, where the alternatives with a
%	singleton source are impossible.
%	
%	\ref{prelim:prop:2homogeneous-pair:3}$\Rightarrow$\ref{prelim:prop:2homogeneous-pair:1}:
%	Let $O$ be the orbit of $\{a,b\}$, which by assumption contains every $2$-set meeting
%	$\{a,b\}$, and let $\{x,y\}$ be disjoint from $\{a,b\}$.
%	Choose $g \in G$ with $g(\{a,b\}) = \{a,x\}$. Then put
%	\[
%	S = 
%	\begin{cases}
%		\{a, g^{-1}(y)\} & \text{if } g(a) = x, \\
%		\{b, g^{-1}(y)\} & \text{if } g(b) = x.	
%	\end{cases}
%	\]
%	As $g$ maps $a$ and $b$ into $\{a,x\}$ and $y \notin \{a,x\}$, the point $g^{-1}(y)$ is
%	neither $a$ nor $b$, so $S$ is a $2$-set meeting $\{a,b\}$ and thus lies in $O$.
%	Now $g(S) = \{x,y\}$ shows $\{x,y\} \in O$, so all $2$-sets form a single orbit.
%\end{proof}

\begin{proposition} \label{prelim:prop:2homogeneous-pair}
	Let $G \le_G \mathcal S_n$ with $n \ge 3$. Then the following are equivalent.
	\begin{enumerate}
		\item \label{prelim:prop:2homogeneous-pair:1} $G$ is $2$-homogeneous.
		
		\item \label{prelim:prop:2homogeneous-pair:2} There exist distinct $a,b \in [n]$ such
		that for all $c,d \in [n]$ with $\{c,d\} \ne \{a,b\}$ there is $g \in G$ with
		$g(\{a,c\}) = \{a,d\}$ or $g(\{b,c\}) = \{b,d\}$.
				
		\item \label{prelim:prop:2homogeneous-pair:4} There exists $a \in [n]$ such that
		\begin{enumerate}
			\item \label{prelim:prop:2homogeneous-pair:4a} for all
			$b,c \in [n] \setminus \{a\}$ there is $g \in G$ with $g(\{a,b\}) = \{a,c\}$, and
			
			\item \label{prelim:prop:2homogeneous-pair:4b} there are
			$b \in [n] \setminus \{a\}$ and $g \in G$ with $a \notin g(\{a,b\})$.
		\end{enumerate}
	\end{enumerate}
\end{proposition}
\begin{proof}
\ref{prelim:prop:2homogeneous-pair:1}$\Rightarrow$\ref{prelim:prop:2homogeneous-pair:2}:
Let $a,b$ be any two distinct points and consider $\{c,d\} \ne \{a,b\}$.
Then $a \notin \{c,d\}$ or $b \notin \{c,d\}$.
As the two conditions in \cref{prelim:prop:2homogeneous-pair:2} arise from one another by
interchanging $a$ and $b$, we may assume $a \notin \{c,d\}$.
Then $\{a,c\}$ and $\{a,d\}$ are $2$-sets and $2$-homogeneity gives
$g(\{a,c\}) = \{a,d\}$ for some $g \in G$.
	
	\ref{prelim:prop:2homogeneous-pair:2}$\Rightarrow$\ref{prelim:prop:2homogeneous-pair:4}:
	Let $a,b$ be as in \cref{prelim:prop:2homogeneous-pair:2} and let $O$ be the orbit of
	$\{a,b\}$.
	Note that $g(\{x,x\})$ is a singleton and hence never a $2$-set.
	For $c \ne a$ applying \cref{prelim:prop:2homogeneous-pair:2} to $b,c$, which is admissible
	as $\{b,c\} \ne \{a,b\}$, thus gives $g(\{a,b\}) = \{a,c\}$, so that $\{a,c\} \in O$.
	Hence any two $2$-sets containing~$a$ lie in $O$ and are therefore images of one another.
	Finally, picking $e \notin \{a,b\}$, possible as $n \ge 3$, and applying
	\ref{prelim:prop:2homogeneous-pair:2} to $a,e$, admissible as $\{a,e\} \ne \{a,b\}$, gives
	$g(\{b,a\}) = \{b,e\}$ with $a \notin \{b,e\}$.

\ref{prelim:prop:2homogeneous-pair:4}$\Rightarrow$\ref{prelim:prop:2homogeneous-pair:1}:
By assumption, there exist $b \in [n] \setminus \{a\}$ and $g \in G$
such that $a \notin g(\{a,b\})$.
Let $O$ be the orbit of $\{a,b\}$ for the action of~$G$ on the $2$-sets.
By assumption, every $2$-set containing~$a$ lies in $O$.
This also implies that every $2$-set containing
a point from $\Orbit{a}{G}$ is in~$O$
since each such $2$-set can be mapped to one containing~$a$.

Put $d = g(a)$, so that $d \ne a$. For $z \ne a$
the $2$-sets $\{a,z\}$ and $\{a,d\}$ lie in~$O$.
So $h(\{a,d\}) = \{a,z\}$ for some $h \in G$.
Whence $z = h(a)$ or $z = h(d) = h(g(a))$
and thus $z \in \Orbit{a}{G}$.
So $\Orbit{a}{G} = [n]$, and by the above every $2$-set lies in~$O$.
\end{proof}

\begin{remark}
	Condition \ref{prelim:prop:2homogeneous-pair:4b} cannot be omitted from
	\cref{prelim:prop:2homogeneous-pair:4} of \cref{prelim:prop:2homogeneous-pair}:
	for $G = \langle (1~2~3) \rangle_G$ on $\{0,1,2,3\}$ and $a = 0$ the three $2$-sets
	containing~$0$ are images of one another, so that
	\ref{prelim:prop:2homogeneous-pair:4a} holds, yet $G$ is not $2$-homogeneous, and not even
	transitive.
\end{remark}


The following is immediate but stated for later reference.

\begin{proposition} \label{syncmax:prop:intransitive-subset}
	Let $G$ be a permutation group on $[n]$ with $n \ge 2$ that is not transitive.
	Then there exists a proper subset $A \subsetneq [n]$ with $|A| \ge \lceil \tfrac{n}{2} \rceil$
	such that $g(A) = A$ for all $g \in G$.
\end{proposition}
\begin{proof}
	Let $O$ be an orbit of $G$. Both $O$ and $[n] \setminus O$ are $G$-invariant, and one of them
	has size at least $\lceil \tfrac{n}{2} \rceil$.
	As $G$ is not transitive, both are proper subsets of $[n]$.
\end{proof}








%
%
%\subsection{}
%
%\begin{comment}
%
%\subsection{Separating sets and idempotents}
%
%In transitive permutation groups we can separate subsets in the following way
%(see~\cite[p.~130]{ArCamSt2017} for related discussion).
%
%\begin{theorem}[\cite{ArCamSt2017,birch2009,neumann75b}]
%	\label{thm:separating_neumann}
%	Let $G \le_G \mathcal S_n$ be transitive and let $A,B \subseteq [n]$ with $|A| \cdot |B| < n$.
%	Then there exists $g \in G$ such that $g(A) \cap B = \emptyset$.
%\end{theorem}
%
%Throughout, for $f \colon [n] \to [n]$ we write $\dupl{f} = \{ y \in [n] \mid |f^{-1}(y)| \ge 2 \}$
%and $\excl{f} = [n] \setminus f([n])$; the \emph{defect} of $f$ is $n - |f([n])|$.
%We call $f$ \emph{one-point singular} if $|\dupl{f}| = 1$, and a \emph{simple idempotent}
%is an idempotent of defect one.
%Note that every transformation of defect one is one-point singular, but not conversely.
%
%We first record that a transitive group together with a transformation of defect one
%generates at least two simple idempotents.
%
%\begin{lemma}
%	\label{sync_words:lem:idempotent}
%	Let $G \le_G \mathcal S_n$ be a transitive permutation group and let $t \colon [n] \to [n]$
%	be of defect one with $t(a) = t(b)$ for distinct $a,b \in [n]$.
%	Then there exists a simple idempotent $h \in \langle G \cup \{t\} \rangle_S$
%	with $h(a) = h(b) = b$.
%\end{lemma}
%\begin{proof}
%	Write $\excl{t} = \{c\}$. Since $G$ is transitive there exists $g \in G$ with $g(c) = a$,
%	and then $\excl{gt} = \{a\}$, i.e.\ $gt([n]) = [n] \setminus \{a\}$.
%	As $gt$ has defect one and $gt(a) = gt(b)$, the restriction of $gt$ to $[n] \setminus \{a\}$
%	is injective, and its image is $gt([n]) = [n] \setminus \{a\}$.
%	Hence $gt$ permutes $[n] \setminus \{a\}$, so there exists $i > 0$ such that
%	$h = (gt)^i$ acts as the identity on $[n] \setminus \{a\}$.
%	In particular $h$ is idempotent of defect one, i.e.\ a simple idempotent, and
%	\[
%	h(a) = (gt)^{i-1}\bigl(gt(a)\bigr) = (gt)^{i-1}\bigl(gt(b)\bigr) = h(b) = b ,
%	\]
%	the last equality because $b \ne a$.
%\end{proof}
%
%Applying \cref{sync_words:lem:idempotent} to the pair $(a,b)$ and to the pair $(b,a)$ shows that
%for $t$ of defect one with $t(a) = t(b)$ the semigroup $\langle G \cup \{t\} \rangle_S$
%contains \emph{both} simple idempotents identifying $a$ and $b$.
%
%\begin{example}
%	Let $g = (0 ~ 1 ~ 2 ~ 3)$ and let $t \colon [4] \to [4]$ be given by
%	$t(0) = 1$, $t(2) = 3$ and $t(1) = t(3) = 2$.
%	Then $gt$ permutes $\{0,2,3\}$ and $g^3 t$ permutes $\{0,1,2\}$.
%	Here $g^3 t$ is already the simple idempotent mapping $3$ to $1$,
%	and $(gt)^2$ is the simple idempotent mapping $1$ to $3$.
%\end{example}
%
%For one-point singular transformations of arbitrary defect we still obtain an idempotent,
%though no longer a simple one; instead, the whole fibre over the duplicated point is preserved.
%
%\begin{lemma} \label{prelim:lem:transitive-grp-one-point-singular}
%	Let $G \le_G \mathcal S_n$ be a transitive permutation group and let $t \colon [n] \to [n]$
%	be one-point singular.
%	Then for each $a \in t^{-1}(\dupl{t})$ there exists an idempotent one-point singular
%	$h \in \langle G \cup \{t\} \rangle_S$ with $\dupl{h} = \{a\}$ and
%	$t^{-1}(\dupl{t}) \subseteq h^{-1}(a)$.
%\end{lemma}
%\begin{proof}
%	Write $\dupl{t} = \{b\}$ and fix $a \in t^{-1}(b)$.
%	Since $G$ is transitive there is $g \in G$ with $g(b) = a$; set $f = gt$.
%	As $g$ is a bijection we get
%	\[
%	\dupl{f} = g(\dupl{t}) = \{a\}, \qquad
%	f(a) = g(t(a)) = g(b) = a, \qquad
%	f^{-1}(a) = t^{-1}(g^{-1}(a)) = t^{-1}(b).
%	\]
%	
%	We claim $\dupl{f^j} = \{a\}$ for all $j \ge 1$.
%	For the inclusion $\subseteq$, let $x \ne y$ with $f^j(x) = f^j(y)$ and let $k \le j$
%	be minimal with $f^{k}(x) = f^{k}(y)$; then $k \ge 1$ and
%	$f^{k-1}(x) \ne f^{k-1}(y)$ are two distinct $f$-preimages of $f^{k}(x)$,
%	whence $f^{k}(x) \in \dupl{f} = \{a\}$.
%	Since $f(a) = a$ this yields $f^{j}(x) = f^{j-k}(a) = a$.
%	For the reverse inclusion, pick $x \ne y$ with $f(x) = f(y) = a$;
%	then $f^{j}(x) = f^{j}(y) = a$, again because $a$ is fixed.
%	
%	Moreover $f^{-1}(a) \subseteq (f^{j})^{-1}(a)$ for every $j \ge 1$: if $f(c) = a$ then
%	$f^{j}(c) = f^{j-1}(a) = a$.
%	
%	By \cref{prelim:lem:finite-semigrp-idempotent-power} there is $i > 0$ such that
%	$h = f^{i}$ is idempotent.
%	By the above, $\dupl{h} = \{a\}$, so $h$ is one-point singular, and
%	$t^{-1}(\dupl{t}) = f^{-1}(a) \subseteq h^{-1}(a)$.
%\end{proof}
%
%\begin{remark}
%	\Cref{sync_words:lem:idempotent} is not a special case of
%	\cref{prelim:lem:transitive-grp-one-point-singular}:
%	the latter applies to $t$ of defect one as well, but the idempotent it produces
%	is only guaranteed to be one-point singular, since preimages may accumulate when passing
%	to powers of $f$. \Cref{sync_words:lem:idempotent} instead chooses $g$ so that $gt$ excludes
%	exactly the point it duplicates, which keeps the defect equal to one along all powers.
%\end{remark}
%
%\subsection{Blocks and point stabilisers}
%
%We next establish a correspondence between the subgroups containing the stabiliser of a point $a$
%and the blocks containing $a$.
%\Crefitem{prelim:thm:blocks-stabs}{prelim:thm:blocks-stabs:2} also appears
%in~\cite[Thm.~1.5A]{dixon1996permutation}.
%
%\begin{theorem}\label{prelim:thm:blocks-stabs}
%	Let $G \le_G \mathcal S_n$ be transitive and $a \in [n]$. Then the following hold.
%	\begin{enumerate}
%		\item \label{prelim:thm:blocks-stabs:2} The maps
%		\[
%		H \mapsto \Orbit{a}{H}
%		\qquad\text{and}\qquad
%		\Delta \mapsto \{ g \in G \mid g(\Delta) = \Delta \}
%		\]
%		are mutually inverse bijections between the subgroups $H$ with $\Stab{a}{G} \le_G H \le_G G$
%		and the blocks $\Delta$ of $G$ with $a \in \Delta$.
%		
%		\item \label{prelim:thm:blocks-stabs:3} For $\Stab{a}{G} \le_G H \le_G G$,
%		$\Delta = \Orbit{a}{H}$ and $g_1,g_2 \in G$ we have
%		\[
%		g_1 H = g_2 H \iff g_1(\Delta) = g_2(\Delta),
%		\]
%		and hence $gH = \{ g' \in G \mid g'(\Delta) = g(\Delta) \}$.
%	\end{enumerate}
%\end{theorem}
%\begin{proof}
%	\ref{prelim:thm:blocks-stabs:2}:
%	Let $\Stab{a}{G} \le_G H \le_G G$ and $\Delta = \Orbit{a}{H}$.
%	We show
%	\begin{equation}\label{prelim:eqn:block-stab-key}
%		\{ g \in G \mid g(\Delta) \cap \Delta \ne \emptyset \} = H .
%	\end{equation}
%	If $g \in H$ then $g(\Delta) = \Delta$, since $\Delta$ is an $H$-orbit; in particular the
%	left-hand side contains $H$.
%	Conversely, suppose $g(h_2(a)) = h_1(a)$ for some $h_1,h_2 \in H$.
%	Then $h_1^{-1} g h_2 \in \Stab{a}{G} \le_G H$ and therefore $g \in h_1 H h_2^{-1} = H$.
%	This proves~\eqref{prelim:eqn:block-stab-key}, and it shows both that $\Delta$ is a block
%	and that the second map returns $H$ when applied to $\Delta$.
%	
%	Now let $\Delta$ be a block with $a \in \Delta$ and set $H = \{ g \in G \mid g(\Delta) = \Delta \}$,
%	which is a subgroup.
%	If $g \in \Stab{a}{G}$ then $a \in g(\Delta) \cap \Delta$, so $g(\Delta) = \Delta$ and $g \in H$;
%	hence $\Stab{a}{G} \le_G H$.
%	The inclusion $\Orbit{a}{H} \subseteq \Delta$ is immediate from the definition of $H$.
%	Conversely let $b \in \Delta$. By transitivity $b = g(a)$ for some $g \in G$, so
%	$b \in g(\Delta) \cap \Delta$ and thus $g(\Delta) = \Delta$, i.e.\ $g \in H$ and $b \in \Orbit{a}{H}$.
%	Hence $\Delta = \Orbit{a}{H}$, so the first map returns $\Delta$, and the two maps are
%	mutually inverse.
%	
%	\ref{prelim:thm:blocks-stabs:3}:
%	By~\eqref{prelim:eqn:block-stab-key} we have $g \in H$ if and only if $g(\Delta) = \Delta$. Thus
%	\[
%	g_1 H = g_2 H
%	\iff g_2^{-1} g_1 \in H
%	\iff g_2^{-1} g_1 (\Delta) = \Delta
%	\iff g_1(\Delta) = g_2(\Delta). \qedhere
%	\]
%\end{proof}
%
%\IN{orbit-stabilizer property}
%\begin{remark}\label{prelim:rem:orbit-stab}
%	Specialising \crefitem{prelim:thm:blocks-stabs}{prelim:thm:blocks-stabs:3} to $H = \Stab{a}{G}$,
%	the left cosets of $\Stab{a}{G}$ are in bijection with the points of $[n]$.
%	As the cosets partition $G$ and all have size $|\Stab{a}{G}|$, we obtain
%	$|G| = n \cdot |\Stab{a}{G}|$, the \emph{orbit-stabilizer property}
%	\cite{cameron_1999,dixon1996permutation,Kurzweil04,wielandt1964finite}.
%	In particular, a transitive $G \le_G \mathcal S_n$ is regular if and only if $|G| = n$.
%\end{remark}
%
%\begin{corollary}
%	\label{sync_max:cor:point_stabs}
%	Let $G \le_G \mathcal S_n$ be transitive, $a \in [n]$ and $\Stab{a}{G} \le_G H \le_G G$.
%	Set $\Delta = \Orbit{a}{H}$. Then:
%	\begin{enumerate}
%		\item \label{sync_max:cor:point_stabs:1} For $g \in G$ we have $g \notin H$ if and only if $a \notin g(\Delta)$.
%		\item \label{sync_max:cor:point_stabs:Htrans} If $H \ne G$ then $H$ does not act transitively.
%	\end{enumerate}
%\end{corollary}
%\begin{proof}
%	\ref{sync_max:cor:point_stabs:1}: Since $\Delta$ is a block containing $a$,
%	the condition $a \in g(\Delta)$ is equivalent to $g(\Delta) \cap \Delta \ne \emptyset$
%	and hence, by~\eqref{prelim:eqn:block-stab-key}, to $g \in H$.
%	
%	\ref{sync_max:cor:point_stabs:Htrans}: By
%	\crefitem{prelim:thm:blocks-stabs}{prelim:thm:blocks-stabs:2} the assignment
%	$H \mapsto \Orbit{a}{H}$ is injective, and $\Orbit{a}{G} = [n]$ by transitivity of $G$.
%	Hence $\Orbit{a}{H} = [n]$ forces $H = G$.
%\end{proof}
%
%\begin{remark}\label{prelim:permgrp:Gactions-rep}
%	By \crefitem{prelim:thm:blocks-stabs}{prelim:thm:blocks-stabs:3} the map $\varphi(gH) = g(\Delta)$
%	is a well-defined bijection from the cosets of $H$ onto the block system generated by
%	$\Delta = \Orbit{a}{H}$, and $\varphi(h \cdot gH) = h(\varphi(gH))$ for $h \in G$.
%	Thus subgroups containing $\Stab{a}{G}$ together with their cosets correspond to blocks
%	containing $a$ together with their block systems, compatibly with the actions by left
%	multiplication and by function application.\footnote{See e.g.~\cite{cameron_1999} for the
%		formalisation via $G$-actions; note that by our conventions permutations act on the left,
%		whereas in~\cite{cameron_1999} they act on the right.}
%	Concretely, for $g \in G$ and $h \in G$,
%	\begin{equation}\label{eqn:congruences-subgroups-blocks}
%		h \in gH \iff h(a) \in g(\Delta) .
%	\end{equation}
%	Indeed, if $h \in gH$ then $h(\Delta) = g(\Delta)$ and $h(a) \in h(\Delta)$;
%	conversely, if $h(a) \in g(\Delta) = \Orbit{a}{gH}$ then $g^{-1}h(a) = h'(a)$ for some $h' \in H$,
%	so $h'^{-1} g^{-1} h \in \Stab{a}{G} \le_G H$ and $h \in gH$.
%	Consequently $G$-congruences on $[n]$ transfer to $G$-congruences on the cosets of $H$.
%	
%	If $G$ is regular, $\varphi$ identifies $G$ with its permutation domain and may be identified
%	with the map $g \mapsto g(a)$.
%	Finally, let $G$ be as in \cref{lem:congruence_on_cyclic_group}.
%	Then $G$ can be identified with the additive group of residues modulo $n$ via
%	$g^{i} \mapsto [i]_{\equiv_n}$.
%	Since every subgroup of that group is the group of multiples of some divisor $r$ of $n$
%	\cite[Sec.~1.4]{Kurzweil04}, \eqref{eqn:congruences-subgroups-blocks} recovers
%	\cref{lem:congruence_on_cyclic_group}.
%\end{remark}
%
%The following is immediate but stated for later reference.
%
%\begin{proposition} \label{syncmax:prop:intransitive-subset}
%	Let $G$ be a permutation group on $[n]$ with $n \ge 2$ that is not transitive.
%	Then there exists a proper subset $A \subsetneq [n]$ with $|A| \ge \lceil \tfrac{n}{2} \rceil$
%	such that $g(A) = A$ for all $g \in G$.
%\end{proposition}
%\begin{proof}
%	Let $O$ be an orbit of $G$. Both $O$ and $[n] \setminus O$ are $G$-invariant, and one of them
%	has size at least $\lceil \tfrac{n}{2} \rceil$.
%	As $G$ is not transitive, both are proper subsets of $[n]$.
%\end{proof}
%
%\end{comment}


\todo[inline]{noch auskommentiertes}

\begin{comment}


\subsection{alt}


Next, we establish a correspondence between certain subgroups containing the point stabilizer of a point $a$ and blocks containing~$a$.
\Crefitem{prelim:thm:blocks-stabs}{prelim:thm:blocks-stabs:2}
also appears in~\cite[Thm.~1.5A]{dixon1996permutation}. 

\begin{theorem}\label{prelim:thm:blocks-stabs}
	Let $G \le_G \mathcal S_n$ be transitive and $a \in [n]$. Then the following is true.	
	
	\begin{enumerate}
		%\item \label{prelim:thm:blocks-stabs:1} $\Orbit{a}{H}$ is a block for $\Stab{a}{G} \le_G H \le_G G$.
		
		\item \label{prelim:thm:blocks-stabs:2} The maps 
		\[
		H \mapsto \Orbit{a}{H}
		\]
		for  $\Stab{a}{G} \le_G H \le_G G$,
		and 
		\[
		\Delta \mapsto \{ g \in G \mid g(\Delta) = \Delta \}
		\]
		for $\Delta$ a block with $a \in \Delta$,
		are mutually inverse mappings between the
		subgroups containing $\Stab{a}{G}$
		and the blocks of $G$ containing $a$, respectively.
		%	that respect set inclusion. \todo{genauer als order morphism formulieren}
		
		\item \label{prelim:thm:blocks-stabs:3} For
		$\Stab{a}{G} \le_G H \le_G G$, $\Delta = \Orbit{a}{H}$,
		and $g_1, g_2 \in G$ we have
		\[
		g_1 H = g_2 H \Leftrightarrow g_1(\Delta)  = g_2(\Delta),
		\]	
		and so $gH = \{ g' \in G \mid g'(\Delta) = g(\Delta)\}$.
	\end{enumerate}
\end{theorem}
\begin{proof}
	\ref{prelim:thm:blocks-stabs:2}: First, let $\Stab{a}{G} \le_G H \le_G G$. We show that $\Orbit{a}{H}$ is a block.
	For $g \in G$, we have
	\begin{align*}
		\Orbit{a}{H} \cap \applyfunc{\Orbit{a}{H}}{g} \ne \emptyset
		& \Leftrightarrow 
		\exists h_1, h_2 \in H \colon \applyfunc{a}{h_1} = \applyfunc{\applyfunc{a}{h_2}}{g} \\
		& \Leftrightarrow 
		a = \applyfunc{\applyfunc{\applyfunc{a}{h_2}}{g}}{h_1^{-1}} \\
		& \Leftrightarrow \concat{\concat{h_2}{g}}{h_1^{-1}} \in G_a.
	\end{align*}
	As $G_a \le_G H$, we find $g \in h_1 H h_2^{-1} = H$.	
	Thus, $g(\Orbit{a}{H}) = \Orbit{a}{H}$
	implies $g \in H$, so that 
	$
	\{ g \in G \mid g(\Orbit{a}{H}) = \Orbit{a}{H} \} \subseteq H.
	$
	Next, clearly we have $H \subseteq \{ g \in G \mid g(\Orbit{a}{H}) = \Orbit{a}{H} \}$
	and since already $\Orbit{a}{H} \cap \applyfunc{\Orbit{a}{H}}{g} \ne \emptyset$
	implies $g \in H$, from this we get $g(\Orbit{a}{H}) = \Orbit{a}{H}$
	so that $\Orbit{a}{H}$ is a block.
	
	%As $H \subseteq \{ g \in G \mid g(\Orbit{a}{H}) = \Orbit{a}{H} \}$, $g \in H$
	%implies $\applyfunc{\Orbit{a}{H}}{g} = \Orbit{a}{H}$.
	%So, $\Orbit{a}{H}$ is a block containing $a$
	%and the second map gives $H$ again. 
	
	Next, let $\Delta$ be a block with $a \in \Delta$.
	Set $H = \{ g \in G \mid g(\Delta) = \Delta \}$,
	which is a subgroup. Consider $g \in \Stab{a}{G}$.
	Then $a \in g(\Delta) \cap \Delta$
	and so $g(\Delta) = \Delta$, i.e., $g \in H$.
	Hence, $G_a \le_G H$.
	By the definition of $H$, we have $\Orbit{a}{H} \subseteq \Delta$.
	Let $b \in \Delta$. As $G$ is transitive,
	we have $b = g(a)$ for some $g \in G$.
	As $\Delta$ is a block, we find $g(\Delta) = \Delta$, i.e., $g \in H$.
	Hence $\Delta \subseteq \Orbit{a}{H}$. Now $\Delta = \Orbit{a}{H}$  means that
	the first map yields $\Delta$ again when applied to the result of the second map.
	So, both maps are mutually inverse.
	
	Thus, we have
	\[
	g' \in gH \Leftrightarrow g'H = gH \Leftrightarrow g'(\Delta) = g(\Delta)
	\]
	and so $gH = \{ g' \in G \mid g'(\Delta) = g(\Delta) \}$.
	
	
%	Suppose there exists a subgroup $H \le G$
%	containing $G_x$. Set $\Delta = \{ \applyfunc{x}{h} \mid h \in H \}$. 
%	
%	Let $x \in [n]$.
%	Let $g \in G$. If $x \in \applyfunc{\Delta}{g}$,
%	then $x = \applyfunc{\applyfunc{x}{h}}{g}$
%	for some $h \in H$.
%	So $hg \in G_x$. As $G_x \le H$, we find $g \in h^{-1}H = H$.
%	Hence for each $g \in G \setminus H$
%	we have $x \notin \applyfunc{\Delta}{g}$.
%	
%	
%	Furthermore, If $H \ne G$, then $H$ cannot act transitive. For if $\applyfunc{x}{g} = \applyfunc{x}{h}$
%	for some $g \in G$, then $x = \applyfunc{\applyfunc{x}{h}}{g^{-1}}$.
%	Hence $hg^{-1} \in G_x \le H$,
%	implying $g^{-1} \in h^{-1}H = H$.
	
	
	
	\ref{prelim:thm:blocks-stabs:3}: Let $\Stab{a}{G} \le_G H \le_G G$ and $\Delta = \Orbit{a}{H}$.
	By \ref{prelim:thm:blocks-stabs:2} for $g \in G$, we have $g \in H$ if and only if $g(\Delta) = \Delta$.
	So
	$
	g_1H = g_2H
	\Leftrightarrow \exists h \in H \colon g_1 = g_2 h 
	\Leftrightarrow g_2^{-1} g_1 \in H 
	\Leftrightarrow g_2^{-1} g_1(\Delta) = \Delta
	\Leftrightarrow g_1(\Delta) = g_2(\Delta). 
	$ 
\end{proof} 

\IN{orbit-stabilizer property}
\begin{remark}\label{prelim:rem:orbit-stab}
	By \crefitem{prelim:thm:blocks-stabs}{prelim:thm:blocks-stabs:2} 
	the blocks $\Delta$ of $G$ with $a \in \Delta$ are in bijective correspondence with the subgroups $H \le_G G$
	containing $\Stab{a}{G}$.
	By \crefitem{prelim:thm:blocks-stabs}{prelim:thm:blocks-stabs:3},
	for $H \le_G G$ containing $\Stab{a}{G}$
	and $\Delta = \Orbit{a}{H}$,
	the left cosets of $H$ in $G$ 
	are in bijective correspondence with the sets $g(\Delta)$, $g \in G$.
	%\footnote{
		%	When considering the sets $g^{-1}(\Delta)$ it is also possible
		%	to establish a correspondence between the action on the right cosets of $H$
		%	by right multiplication and the action of the inverse elements on the
		%	sets $g^{-1}(\Delta)$ by multiplication on the left, but as we do not need
		%	this result here we do not make this more formal.}
	For the special case of transitive $G \le_G\mathcal S_n$ and
	$H = \Stab{a}{G}$, we obtain that the left cosets of $\Stab{a}{G}$
	are in bijection with the points from $[n]$.
	As the left cosets partition $G$ and all have the same size~$|\Stab{a}{G}|$,
	we get $|G| = n |\Stab{a}{G}|$, which
	is known as the \emph{orbit-stabilizer property} in the literature~\cite{cameron_1999,dixon1996permutation,Kurzweil04,wielandt1964finite}.
	This implies that a transitive group $G \le \mathcal S_n$
	is regular if and only if $|G| = n$.
	We note the following corollary to \cref{prelim:thm:blocks-stabs}.
\end{remark}

\begin{corollary}
	\label{sync_max:cor:point_stabs}
	Let $G \le_G \mathcal S_n$ be transitive, 
	$a \in [n]$ and $\Stab{a}{G} \le_G H \le_G G$.
	Set $\Delta = \Orbit{a}{H}$.
	Then the following properties are true:
	\begin{enumerate}
		\item \label{sync_max:cor:point_stabs:1}  For each $g \in G$ we have $g \in G \setminus H$ if and only if $a \notin \applyfunc{\Delta}{g}$.
		
		\item \label{sync_max:cor:point_stabs:Htrans} If $H \ne G$, then $H$ does not act transitively. 
	\end{enumerate}
\end{corollary}
\begin{proof}
	\ref{sync_max:cor:point_stabs:1}: Since $\Delta$ is a block by \cref{prelim:thm:blocks-stabs}, $a \in g(\Delta)$
	is equivalent to $g(\Delta) = \Delta$ for $g \in G$.
	Since $H = \{ g\in G \mid g(\Delta) = \Delta \}$
	by \crefitem{prelim:thm:blocks-stabs}{prelim:thm:blocks-stabs:2},
	$g \in G \setminus H$ if and only if $a \notin g(\Delta)$.
	
	\ref{sync_max:cor:point_stabs:Htrans}: Let $a \in [n]$. Now, $G$ is transitive
	and,  by \crefitem{prelim:thm:blocks-stabs}{prelim:thm:blocks-stabs:2}, $\Orbit{a}{G} = [n]$ gives
	that  for $H \le_G G$, we have $\Orbit{a}{H} = [n]$
	if and only if $H = G$. Thus we obtain the claim.
\end{proof}

%
% kommutative diagramme zeichnen und G-actions einführen?
\begin{remark}\label{prelim:permgrp:Gactions-rep}
	By \crefitem{prelim:thm:blocks-stabs}{prelim:thm:blocks-stabs:3}
	the map $\varphi$ from the cosets of $H$
	to the system of blocks for $\Delta = \Orbit{a}{H}$
	given by $\varphi(gH) = g(\Delta)$
	is a well-defined bijection such
	that $\varphi(h(gH)) = h(\varphi(gH))$.
	Thus, subgroups and their cosets correspond
	to blocks containing~$a$
	and their systems of blocks
	with the respective actions by left multiplication and function application.\footnote{We refer to the literature as e.g.~\cite{cameron_1999} 
		for details on how these actions are defined via so-called $G$-actions, but note that here
		by our conventions on function applications the permutations act from the left whereas in~\cite{cameron_1999} 
		they act on the right.}	
	Furthermore, for $c,d \in [n]$ and $h_1, h_2 \in gH$ with $h_1(a) = c$ and $h_2(a) = d$ we have
	\begin{equation}\label{eqn:congruences-subgroups-blocks}
		h_1, h_2 \in gH \Leftrightarrow c,d \in g(\Delta),
	\end{equation}
	and so $G$-congruences transfer to $G$-congruences for the action on the cosets of~$H$.
	
	If~$G$ is regular, the~$\varphi$
	gives a correspondence between the elements in~$G$ and its permutation domain
	and the map~$\varphi$ can be identified with the map from~$G$
	to~$[n]$ that maps $g \in G$ to $g(a)$.	
	
	Now, consider a group~$G$ as in \cref{lem:congruence_on_cyclic_group}.
	Then~$G$ can be identified with the abstract group
	given by addition on the residues modulo~$n$
	by identifying $g^i$ with $[i]_{\equiv_n}$
	(we note that also the abstract group of residues modulo~$n$
	can be identified with the regular permutation
	group as in  \cref{lem:congruence_on_cyclic_group}
	by letting~$G$ act on itself via multiplication from left).
	In this case, since the subgroups of the group of residues
	can be identified with groups on the residues modulo~$r$
	for some divisor~$r$ of~$n$, see~\cite[Sec.~1.4]{Kurzweil04},
	\cref{eqn:congruences-subgroups-blocks}
	also yields \cref{lem:congruence_on_cyclic_group}.
\end{remark}

The following is easy to see, but stated for reference since we need it later.

\begin{proposition} \label{syncmax:prop:intransitive-subset}
	Suppose $G$ is a permutation group on $[n]$ with $n \ge 2$ that is not transitive.
	Then there exists a proper subset $A \subseteq [n]$
	with $|A| \ge \lceil \frac{n}{2} \rceil$
	such that $g(A) = A$ for all $g \in G$.
\end{proposition}
\begin{proof}
	Consider some orbit of~$G$.
	Then this orbit or its complement 
	has size at least $\lceil \frac{n}{2} \rceil$
	and both sets fulfill the statement.
	Since~$G$ is not transitive, both sets
	are proper subsets of~$[n]$.
\end{proof}



\end{comment}







%Lastly, we give results relating mapping properties between $2$-sets to orbit sizes
%and transitivity.
%
%The following result relating $k$-homogeneity to $k$-transitivity
%is due to Livingstone and Wagner~\cite{Livingstone1965}.
%
%\begin{theorem}[Livingstone and Wagner~\cite{Livingstone1965}]
%	\label{prelim:thm:livingstone}
%	Suppose $G \le_G \mathcal S_n$ is a $k$-homogeneous permutation group with $k \ge 2$ and $2k \le n$.
%	Then $G$ is $(k-1)$-transitive (and so also $(k-1)$-homogeneous)
%	and, for $k \ge 5$, $G$ is $k$-transitive.	
%\end{theorem}	
%
%\begin{lemma} \label{prelim:permgrp:lem:transitive-abac}
%	Let $G \le_G \mathcal S_n$ and $a \in [n]$. Then the following are equivalent.
%	\begin{enumerate}
%		\item \label{prelim:permgrp:lem:transitive-abac:1} $G$ is transitive.
%		
%		\item \label{prelim:permgrp:lem:transitive-abac:2} For distinct $c,d \in [n]$
%		we have $\applyfunc{\{a,b\}}{g} = \{c,d\}$
%		for some $b \in [n]$ and $g \in G$
%		with $g(a) \ne a$.
%		
%		\item \label{prelim:permgrp:lem:transitive-abac:3} For $b \in [n] \setminus\{a\}$
%		we have $g(\{a,c\}) = \{a,b\}$ for some $c \in [n]$ and $g \in G$ with $g(a) \ne a$.
%	\end{enumerate}
%\end{lemma}
%\begin{proof}
%	$\ref{prelim:permgrp:lem:transitive-abac:1}\Rightarrow\ref{prelim:permgrp:lem:transitive-abac:2}$: Without loss of generality, let $c \ne a$.
%	Since~$G$ is transitive, we have $g(a) = c$ for some $g \in G$.
%	Then $g(\{a,b\}) = \{c,d\}$ where $b = g^{-1}(d)$ and $g(a) \ne a$.
%	
%	$\ref{prelim:permgrp:lem:transitive-abac:2}\Rightarrow\ref{prelim:permgrp:lem:transitive-abac:3}$: This follows from \cref{prelim:permgrp:lem:transitive-abac:1} by considering $c = a$ and $d = b$.
%	
%	$\ref{prelim:permgrp:lem:transitive-abac:3}\Rightarrow\ref{prelim:permgrp:lem:transitive-abac:1}$: Let $b \ne a$.
%	Then $g(\{a,c\}) = \{a,b\}$ for some $c \in [n]$
%	and $g\in G$ with $g(a) \ne a$. The latter condition gives $g(a) = b$.
%\end{proof}
%
%\begin{lemma}\label{prelim:lem:2sets-transitive}
%	Let $G \le_G \mathcal S_n$ and $a \in [n]$.
%	Then we have $|\Orbit{a}{G}| \ge n - 1$
%	if and only if for any $2$-set $\{c,d\} \subseteq [n]$
%	we have $\applyfunc{\{a,b\}}{g} = \{c,d\}$
%	for some $g \in G$ and $b \in [n]$. 
%\end{lemma}
%\begin{proof}
%	First, suppose $|\Orbit{a}{G}| \ge n - 1$.
%	Let $\{c,d\} \subseteq [n]$ be a $2$-set.
%	We have $\{c,d\} \cap \Orbit{a}{G} \ne \emptyset$.
%	Without loss of generality, let $c \in \Orbit{a}{G}$.
%	Then $c = g(a)$ for some $g \in G$
%	and with $b = g^{-1}(d)$
%	we find $g(\{a,b\}) = \{c,d\}$.
%	
%	Next, suppose $|\Orbit{a}{G}| \le n - 2$. Then 
%	there exist distinct $c,d \in [n]$
%	with $c,d \notin \Orbit{a}{G}$.
%	Then $g(a) \notin \{c,d\}$ for all $g \in G$.
%	This implies $g(\{a,b\}) \ne \{c,d\}$
%	for all $g \in G$ and $b \in [n]$,
%	since otherwise $g(a) \notin \{c,d\}$
%	for some $g \in G$. Thus, the $2$-set mapping condition does not hold.
%\end{proof}
%
%Note that in \cref{prelim:lem:2sets-transitive} we have $a \ne b$.
%
%\begin{corollary} \label{prelim:lem:2sets-and-botha-transitive}
%	Let $G \le_G \mathcal S_n$ and $a \in [n]$.
%	Suppose $|\Orbit{a}{G}| \ge n - 1$
%	and for $b \in [n]$ with $b \ne a$
%	we have $g(\{a,c\}) = \{a,b\}$
%	for some $g \in G$ and $c \in [n]$ with $c \ne b$.
%	Then~$G$ is transitive.	
%\end{corollary}
%\begin{proof}
%	We have $[n] = \Orbit{a}{G} \cup \{b\}$
%	for some $b \in [n]$.	
%	We can assume $a \ne b$.
%	By assumption, $g(\{a, c\}) = \{a,b\}$
%	for some $g \in G$ and $c \in [n]$ 
%	with $c \ne b$.
%	This gives $g(b) \ne b$ and thus $g(b) \in \Orbit{a}{G}$.
%	Hence $b \in \Orbit{a}{G}$
%	and so~$G$ is transitive.\footnote{For an alternative argument, 
%		consider some $b \in [n]$ with $a \ne b$.
%		By assumption $g(\{a, c\}) = \{a,b\}$ for some $c \in [n]$
%		with $c \ne b$.
%		For $g(a) \ne a$ we have $g(a) = b$.
%		Otherwise, $g(a) = a$ gives $g(c) = b$
%		and since $\{b,c\}$ is a $2$-set, with \cref{prelim:lem:2sets-transitive}, 
%		we have $h(\{a,d\}) = \{b,c\}$ for some $d \in [n]$ and $h \in G$.
%		Then $h(a) = b$ or, for $h(a) = c$, we get $g(h(a)) = b$.
%		Thus, $a$ can be mapped to any other point~$b \in [n]$
%		and so $G$ is transitive.}
%	%Thus $c \in \Orbit{a}{G}$. Since $g(a) = b$ or $g(c) = b$, we find $b \in \Orbit{a}{G}$ and so~$G$ is transitive.
%\end{proof}
%
%\begin{corollary}
%	Let $G \le_G \mathcal S_n$ and $a \in [n]$.
%	Suppose for any $2$-set $\{c,d\} \subseteq [n]$
%	we have  $\applyfunc{\{a,b\}}{g} = \{c,d\}$
%	for some $g \in G$ and $b \in [n]$ with
%	$b \ne c$. Then~$G$ is transitive.
%\end{corollary}
%\begin{proof}
%	By \cref{prelim:lem:2sets-transitive} 
%	we have $|\Orbit{a}{G}| \ge n - 1$.
%	By assumption, for $b \ne a$
%	we have $g(\{a,c\}) = \{a,b\}$
%	for some $c \ne b$. Thus, with \cref{prelim:lem:2sets-and-botha-transitive}, $G$ is transitive.
%\end{proof}
%
%\begin{corollary}
%	Let $G \le_G \mathcal S_n$ and $a \in [n]$.
%	Suppose for any $2$-set $\{c,d\} \subseteq [n]$
%	we have $\applyfunc{\{a,b\}}{g} = \{c,d\}$
%	for some $b \in [n]$ and $g \in G$ with $g(b) \ne b$.
%	Then~$G$ is transitive.	
%\end{corollary}
%\begin{proof}
%	By \cref{prelim:lem:2sets-transitive}, we have $[n] = \Orbit{a}{G} \cup \{b\}$
%	for some $b \in [n]$.
%	We can assume $a \ne b$.
%	By assumption, $g(\{a,c\}) = \{a,b\}$ for some $c \in[n]$
%	and $g\in G$ with $g(c) \ne c$.
%	Then $g(a) = b$
%	or $g(c) = b$ where $c \in \Orbit{a}{G}$.
%	Both cases give $b \in \Orbit{a}{G}$
%	and so~$G$ is transitive.\footnote{Alternatively, let $b \in [n]$ be with $a \ne b$. 
%		By assumption, there exist $c \in [n]$ and $g \in G$
%		such that $g(\{a,c\}) = \{b,a\}$ and $g(c) \ne c$.
%		Suppose $g(a) \ne b$. Then $g(a) = a$ and so $g(c) = b$.
%		Since $g(c) \ne c$, we have $c \ne b$.
%		By assumption, there exist $d \in [n]$ and $h \in G$
%		such that $h(\{a,d\}) = \{b,c\}$.
%		For $h(a) \ne b$, we have $h(a) = c$.
%		Then $g(h(a)) = g(c) = b$.
%		Thus, we can map~$a$ to any other point and so~$G$ is transitive.}
%\end{proof}
%
%\begin{lemma} \label{prelim:permgrp:lem:from-ab-or}
%	Let $G \le_G \mathcal S_n$
%	and $a, b\in [n]$ be distinct.
%	Suppose for $c \in [n]$
%	we have  $g(\{a,b\}) = \{a,c\}$ or $g(\{a,b\}) = \{b,c\}$ 
%	for some $g \in G$.
%	Then $[n] = \Orbit{a}{G} \cup \Orbit{b}{G}$.
%\end{lemma}
%\begin{proof}
%	By assumption, $c \in \Orbit{a}{G} \cup \Orbit{b}{G}$ for $c \in [n]$.
%	Thus $[n] = \Orbit{a}{G} \cup \Orbit{b}{G}$.
%\end{proof}
%
%\begin{corollary} \label{prelim:permgrp:cor:from-ab-or}
%	Let $G \le_G \mathcal S_n$ and $a,b \in [n]$ be distinct.
%	Suppose for $c \in [n]$ with $c \ne a$
%	we have $g(\{a,b\}) = \{a,c\}$ for some $g \in G$.
%	Then $|\Orbit{b}{G}| \ge n - 1$.
%\end{corollary}
%\begin{proof}
%	Let $c \in [n]$.
%	For $c \ne a$ we have $g(\{a,b\}) = \{a,c\}$ for some $g \in G$,
%	and for $c = a$, the identity maps $\{a,b\}$ to itself.
%	Thus, by \cref{prelim:permgrp:lem:from-ab-or} ,
%	we find $[n] = \Orbit{a}{G} \cup \Orbit{b}{G}$.
%	Suppose $c \in \Orbit{a}{G}$ for some $c \ne a$.
%	Then $g(\{a,b\}) = \{a,c\}$ for some $g \in G$.
%	Thus $g(b) \in \{a,c\} \subseteq \Orbit{a}{G}$,
%	which gives $b \in \Orbit{a}{G}$
%	and so $[n] = \Orbit{a}{G}$. This gives the claim.	
%\end{proof}
%
%\begin{corollary}
%	Let $G \le_G \mathcal S_n$ and $a \in [n]$.
%	Suppose for all $b,c \in [n] \setminus \{a\}$
%	we have $g(\{a, b\}) = \{a,c\}$ for some $g \in G$.
%	Then $|\Orbit{b}{G}| \ge n - 1$ for some $b \in [n] \setminus \{a\}$.
%\end{corollary}
%\begin{proof}
%	Fix some $b \in [n] \setminus\{a\}$.
%	By \cref{prelim:permgrp:cor:from-ab-or} 
%	we find $|\Orbit{b}{G}| \ge n - 1$.
%\end{proof}
%
%\begin{corollary} \label{prelim:permgrp:cor:from-ab-or2}
%	Let $G \le_G \mathcal S_n$ with $n \ge 3$ and $a,b \in [n]$ be distinct.
%	Suppose for $c,d \in [n]$ with $\{c,d\} \ne \{a,b\}$  
%	we have $g(\{a,c\}) = \{a, d\}$
%	or $g(\{b,c\}) = \{b,d\}$ for some $g \in G$.
%	Then $G$ is transitive.
%\end{corollary}
%\begin{proof}
%	Let $c \notin \{a,b\}$. Then $\{b,c\} \ne \{a,b\}$.
%	By assumption, $g(\{a,b\}) = \{a,c\}$ for some $g \in G$.
%	Thus $|\Orbit{b}{G}| \ge n - 1$
%	by \cref{prelim:permgrp:cor:from-ab-or}.
%	Similarly, by considering $\{a,c\}$,
%	we find  $|\Orbit{a}{G}| \ge n - 1$.
%	Since $n \ge 3$, this yields that~$G$ is transitive.\footnote{Alternatively, 
%		let $c \in [n]$ with $c \notin \{a,b\}$.
%		Since $\{a,c\} \ne \{a,b\}$
%		we have $g(\{a\}) = \{a, c\}$
%		or $g(\{b, a\}) = \{b,c\}$
%		for some $g \in G$.
%		Now, since $c \ne a$,
%		$g(\{a\}) = \{a,c\}$ cannot be the case.
%		If $g(b) = c$, then $g(a) = b$
%		and we find $g(g(a)) = c$.
%		Otherwise, $g(b) = b$ and so $g(a) = c$.
%		Thus we find $c \in \Orbit{a}{G}$.
%		Similarly $c \in \Orbit{b}{G}$.
%		This gives $[n] = \Orbit{a}{G}$.}
%\end{proof}
%
%The argument used in the proof of the following lemma was given in~\cite[p.~533]{mcalister1998semigroups}
%to deduce that a $2$-homogeneous group is transitive.
%
%\begin{lemma} \label{prelim:lem:2sets-transitive2}
%	Let $G \le_G \mathcal S_n$ with $n \ge 3$ and $a \in [n]$.
%	Suppose for any $c \in [n]$ there exists $b \in [n]$
%	such that $g(\{a,b\}) = \{c,b\}$. 
%	Then $G$ is transitive.
%\end{lemma} 
%\begin{proof}
%	Let $c \in [n]$.
%	Then $g(\{a,b\}) = \{c,b\}$ for some $b \in [n]$ and $g \in G$.
%	For $g(a) = b$, we have $g(b) = c$ so that $g(g(a)) = c$.
%	Otherwise, $g(a) = c$.
%\end{proof}
%
%\begin{remark}
%	Transitivity does not imply the condition from \cref{prelim:lem:2sets-transitive2}.
%	For example, $G = \langle \{(0 ~ 1 ~ \cdots ~ n - 1)\} \rangle_G$ does not fulfill it.
%\end{remark}

\section{Permutation Groups and Abstract Groups}
\label{prelim:sec:perm_grp_abstract_grp}

%!TEX root = ../../springer-mono-format/book.tex
%!TEX root = ../../thesis-trier.tex

\IN{abstract group}
Often groups are given by some representation. For example, by sets of permutations as permutation groups. We call a group an \emph{abstract group}\footnote{In the same vein, we can talk about abstract monoids or abstract semigroups when solely referring to their axiomatic properties.}	
if we want to emphasize that 
we do not assume some internal structure or a \mbox{(re-)}~presentation beyond the axiomatic definition of a group.

\IN{left regular (Cayley) representation}
\IN{diagonal action}
Now, let $(G, \cdot, ^{-1}, 1)$ be an abstract group. Next, we describe one possibility on
how to represent $G$ as a permutation group. For each $g \in G$
let $\lambda_g \colon G \to G$ be the permutation
defined by $\lambda_g(h) = gh$ for $h \in G$.
The map $\Pi \colon G \to \mathcal S_G$ given by $\Pi(g) = \lambda_g$
is injective, and we have $\Pi(gg') = \lambda_{gg'} = \concat{\lambda_{g'}}{\lambda_g}$, $\Pi(1) = \operatorname{id}_G$,
and $\Pi(g^{-1}) = \lambda_{g^{-1}} = \lambda_g^{-1}$ for $g,g' \in G$, i.e., it is a group morphism.
This yields that $G$
can be uniquely identified with the permutation group $\{ \lambda_g \mid g \in G\} \le_G \mathcal S_G$,
giving a permutation representation of $G$ that is called 
the \emph{left regular (Cayley) representation}.
This is not the only way 
to represent an abstract group as a permutation group, 
the interested reader can find more information in the discussion (\cref{prelim:permgrp:sec:discussion}).

\begin{remark}
	Every regular permutation group arises in the form of the 
	left regular representation on itself
	in the following way.
	Let $G \le \mathcal S_n$ be a regular permutation group
	and $a \in [n]$.
	The map $\varphi \colon [n] \to G$ given by $\varphi(b) = g$ for the unique $g \in G$
	with $\applyfunc{a}{g} = b$ is a bijection.
	Since for $h,g \in G$ and $b \in [n]$ with $g(a) = b$ we have $(h \circ g)(a) = h(g(a)) = h(b)$,
	we find
	\[
	\applyfunc{\applyfunc{b}{h}}{\varphi} = h \circ \applyfunc{b}{\varphi} = \applyfunc{\applyfunc{b}{\varphi}}{\lambda_h}
	\] 
	where $\lambda_h \colon G \to G$ is the left regular representation introduced above.
	Thus $h = \varphi^{-1} \lambda_h \varphi$
	and $h$ maps $b$ to $c$ for $b,c \in [n]$
	if and only if $\lambda_h$ maps 
	$\applyfunc{b}{\varphi}$ to $\applyfunc{c}{\varphi}$.
	Hence the map $\varphi$ identifies the permutation domains $[n]$ and~$G$ in a way
	so that~$G$ and~$\overline{G}$ act in the same way on each.
	As an example, applying this construction to the right regular representation $\overline{G} = \{ \rho_g \mid g \in G \}$, which is a subgroup of $\mathcal S_G$, and some $g \in G$,
	we obtain  $\varphi \colon G \to \overline{G}$
	such that for $h \in G$  
	we have $\varphi(h) = \rho_{h^{-1}g}$
	since $\rho_{h^{-1}g}(g) = g(h^{-1}g)^{-1} = h$.
\end{remark}

\IN{diagonal action}
Next, we show how to obtain a representation of $G \times G$
as a subgroup of $\mathcal S_G$ for abstract groups~$G$
with trivial center, called the \emph{diagonal action (on $G$)}.\footnote{This is actually a more general
construction related to so-called group actions, see~\cite[Example~2.6]{ArCamSt2017} or~\cite[Exercise~1.6]{cameron_1999}.} 

\label{def:diagonal-action}
\begin{proposition}
\label{prelim:prop:diagonal_action}
Let $G$ be an abstract group with identity $1 \in G$ and $Z(G) = \{1\}$.
For $x,y \in G$ let $\delta_{x,y} \colon G \to G$
be given by $\delta_{x,y}(g) = xgy^{-1}$.
	
\begin{enumerate}
\item \label{prelim:prop:diagonal_action:1} The map $(x,y) \mapsto \delta_{x,y}$, $x,y \in G$, 
 is an injective morphism from $G\times G$
 into $\mathcal~S_G$.

\item \label{prelim:prop:diagonal_action:2} $G^D = \{ \delta_{x,y} \mid (x,y) \in G\}$ is a transitive permutation group (the representation of $G \times G$ by the diagonal action) on $G$ that is primitive if and only if $G$ is simple.
\end{enumerate}
\end{proposition}
\begin{proof}
\ref{prelim:prop:diagonal_action:1} 
 Let $x_1 x_2, y_1, y_2 \in G$.
 As 
 \[
 \delta_{x_1 x_2, y_1 y_2}(g) = (x_1 x_2) g (y_1 y_2)^{-1} = x_1( x_2 g y_2^{-1}) y_1^{-1} = ( \delta_{x_1,y_1} \delta_{x_2,y_2} )(g),
 \]
 for $g \in G$ we obtain a morphism.
 For injectivity, suppose $\delta_{x_1,y_1} = \delta_{x_2, y_2}$.
 Then $x_2^{-1} x_1 g y_1^{-1} y_2 = g$
for all $g\in G$. For $g = 1$,
we get $x_2^{-1} x_1 y_1^{-1} y_2 = 1$,
which yields $x_2^{-1} x_1 = y_2^{-1} y_1$.
Now, $x_2^{-1} x_1 g (x_2^{-1} x_1)^{-1} = g$ for all $g \in G$
if and only if $x_2^{-1} x_1 \in Z(G)$.
As $Z(G) = \{1\}$, $x_1 = x_2$
and $y_1 = y_2$.
	
\ref{prelim:prop:diagonal_action:2} As $\applyfunc{g}{\delta_{hg^{-1},1}} = hg^{-1}g = h$
for $g,h \in G$, we find that $G^D$ is transitive.
Next, we show that blocks of $G^D$ containing $1$
are precisely the normal subgroups of~$G$, which gives the claim.
Let $\Delta \subseteq G$ be a block of $G^D$ with $1 \in \Delta$ and $x,y \in \Delta$.
As $x = xyy^{-1}$ we have $x \in \Delta \cap x\Delta y^{-1}$.
So $x\Delta y^{-1} = \Delta$.
As $1 \in \Delta$ this yields $xy^{-1} \in \Delta$.
Hence $\Delta$ is a subgroup of $G$.
Now, let $g \in G$.
Then $1 \in \Delta \cap g\Delta g^{-1}$.
So $\Delta = g\Delta g^{-1}$
and $\Delta$ is a normal subgroup.
Conversely, let $N$ be a normal subgroup of $G$.
Suppose $x,y \in G$. We have $\applyfunc{N}{\delta_{x,y}} = 
x N y^{-1} = x y^{-1} y N y^{-1} = x y^{-1} N$, which is a left coset.
Since these cosets partition $G$, $\applyfunc{N}{\delta_{x,y}} \cap N \ne \emptyset$
implies $\applyfunc{N}{\delta_{x,y}} = N$.
So,~$N$ is a block of~$G^D$.
\end{proof}

\IN{right regular (Cayley) representation}
Next, more generally, suppose we have given a subgroup $H$ of~$G$ such that
\footnote{In general,
	$\bigcap_{g \in G} gHg^{-1}$
	is called the \emph{core} of $G$
	and~$H$ is called \emph{core-free} if \cref{prelim:eqn:core} holds.}	
\begin{equation}\label{prelim:eqn:core}
	\bigcap_{g \in G} gHg^{-1} = \{1\}
\end{equation}
Then we can construct a
permutation group on the left cosets $X = \{ hH \mid h \in G\}$
by $\overline{G} = \{ \lambda_g \mid g \in G \}$
where $\lambda_g \colon X \to X$ is a permutation
on $X$ given by $\applyfunc{hH}{\lambda_g} = (gh)H$.
\cref{prelim:eqn:core} ensures that $\lambda_{g} \ne \lambda_{g'}$
if and only if $g \ne g'$.
Further, we have $\lambda_{gg'} = \concat{\lambda_{g'}}{\lambda_{g}}$
for $g, g' \in G$. So, $g \mapsto \lambda_g$ 
is an isomorphism of abstract groups between $G$ and $\overline G$.
In this way, we have realized~$G$ as a permutation group on $X$.
For $H = \{1\}$ (in this case we identify~$X$ and~$G$)
we obtain the left regular representation. 
There is also a corresponding construction on the right cosets 
by considering the mappings $\rho_g \colon Y \to Y$
given by $\applyfunc{Hh}{\rho_g} = H(hg^{-1})$
for $Y = \{ Hh \mid h \in H\}$. The inversion entered by our choice
on the order of function composition,
so that the equation $\rho_{gg'} = \concat{\rho_{g'}}{\rho_{g}}$
is valid (see the related discussion in \cref{prelim:sec:remarks_func_comp}).\footnote{Note that when 
	the convention is adopted that in a function composition the left-most
	function is applied first, we have to use an inversion for the left regular action
	instead, see, for example,~\cite{cameron_1999}.}
The permutation representation obtained here
for $H = \{1\}$ is also called the \emph{right regular (Cayley) representation}.

By considering those subgroups~$H$ that fulfill \cref{prelim:eqn:core} 
we, in fact, obtain all permutation representations of~$G$.
This follows since~$H$ precisely equals the point stabilizer
of the permutation group constructed above, and
that, by \cref{prelim:thm:blocks-stabs}, every permutation group
can be realized as a permutation group on the left cosets for some 
point stabilizer subgroup.
A form of these representations without \cref{prelim:eqn:core} seems to have been first observed by Burnside in~\cite{Burnside1901_representation_permutation}, where he also refers to Dyck (without citing a source).

\section{Permutation Groups and Non-Permutations}
\label{prelim:sec:grp-of-units}

%!TEX root = ./../../thesis-trier.tex

Here, we collect results about transformation semigroups generated by a permutation group~$G$ and a set $T \subseteq \mathcal T_n$ of non-permutations.

\begin{proposition}
	\label{prelim:prop:group_of_units}
	Let $G \le_G \mathcal S_n$ and $T \subseteq \mathcal T_n \setminus S_n$. 
	Then $G$ is the group of units of $\langle G \cup T \rangle_S$.
\end{proposition}
\begin{proof}
		By \crefitem{prelim:prop:rank-defect}{prelim:prop:item:rank-eqn},
		for $t \in \mathcal T_n \setminus \mathcal S_n$
		or $h \in \mathcal T_n$
		we have $\concat{t}{h}, \concat{h}{t}  \in \mathcal T_n \setminus S_n$,
		which gives the claim.
\end{proof}

By successively using the formula $gt = t^g g$ for a transformation $t$
and permutation $g$, we obtain the next statement.

\begin{proposition} 
	\label{sync_max:prop:factorization_KG}
	Let $G \le_G \mathcal S_n$ be a permutation group and $T \subseteq \mathcal T_n$. 
	Set $S = \langle G \cup T \rangle_S$. Let~$K = \{ t^g \mid t \in T, g \in G \}$ be the set of conjugates of elements from~$T$.
	Then $S = \langle K \rangle_M G = G \langle K \rangle_M$.
\end{proposition}
\begin{proof}
	As $K \subseteq S$ and $G \subseteq S$,
	we have $\langle K \rangle_M \subseteq S$
	and so also $\langle K \rangle_M G \subseteq S$.
	Next, suppose $s \in S$. Then we have $s = g_0 t_1 g_1 t_2 g_2 \cdots t_{n-1} g_{n-1} t_n g_n$
	and 
	\begin{equation} \label{sync_max:prop:factorization_KG:eqn}
		g_0 t_1 g_1 t_2 g_2 \cdots t_{n-1} g_{n-1} t_n g_n 
		= t_1^{g_0} t_2^{g_0 g_1} \cdots t_{n-1}^{g_0 g_1 \cdots g_{n-2}} t_n^{g_0 g_1 \cdots g_{n-1}} g_0 g_1 \cdots g_n.
	\end{equation}	
	The above formula can be seen by induction on the number of~$t_i$'s
	in the product. If~$s$ is a product of permutations, then 
	the formula clearly holds.
	Otherwise,
	we have 
	\[ 
	g_0 t_1 g_1 t_2 g_2 \cdots t_{n-1} g_{n-1} t_n g_n 
	= t_1^{g_0} g_0 g_1 t_2 g_2 \cdots t_{n-1} g_{n-1} t_n g_n
	\]
	and, by induction hypothesis,
	\[
	g_0 g_1 t_1 g_2 \cdots t_{n-1} g_{n-1} t_n g_n = t_1^{g_0 g_1} \cdots t_{n-1}^{g_0 g_1 \cdots g_{n-2}} t_n^{g_0 g_1 \cdots g_{n-1}} g_0 g_1 \cdots g_n,
	\]
	which yields \cref{sync_max:prop:factorization_KG:eqn}.
	Since $t_1^{g_0} t_2^{g_0 g_1} \cdots t_{n-1}^{g_0 g_1 \cdots g_{n-2}} t_n^{g_0 g_1 \cdots g_{n-1}} \in \langle K \rangle_M$
	and $g_1 g_2 \cdots g_n \in G$
	we obtain $S \subseteq \langle K \rangle_M G$. So $S = \langle K \rangle_M G$.
	Similarly, we can show
	\[
	 g_0 t_1 g_1 t_2 g_2 \cdots t_{n-1} g_{n-1} t_n g_n 
	 = g_0 g_1 \cdots g_n t_1^{(g_1 g_2 \cdots g_n)^{-1}} t_2^{(g_2 \cdots g_n)^{-1}} \cdots t_{n-1}^{(g_{n-1} g_n)^{-1}} t_n^{g_n^{-1}},
	\]
	which then yields $S = G \langle K \rangle_M$.		
\end{proof}

The following corollary was noted in~\cite[Sec.~2]{DBLP:journals/tcs/AraujoBC13}).

\begin{corollary}
	Suppose $G \le_G \mathcal S_n$ is a permutation group and $t \in \mathcal T_n$.
	If $\langle G \cup \{t\} \rangle_S$
	contains a constant, then some constant %in $\langle G \cup \{t\} \rangle_S$
	is a composition of conjugates of~$t$.
\end{corollary}

Next, we note a special case when $t$ is idempotent.

\begin{corollary}
	\label{sync_max:cor:factorization_EG}
	Suppose $G \le_G \mathcal S_n$ is a permutation group and $t \in \mathcal T_n$ is an idempotent.
	Set $S = \langle G \cup \{t\} \rangle_S$. Let~$E$ be the set of idempotents of $S$.
	Then
	$S = \langle E \rangle_M G$.	
\end{corollary}
\begin{proof}
	Let $K = \{ t^g \mid g \in G\}$.
	By Equation~\eqref{prelim:conj} the mappings in $K$ are idempotents,
	so $K \subseteq E$.
	With Proposition~\ref{sync_max:prop:factorization_KG} 
	we find $S = \langle K \rangle_M G \subseteq \langle E \rangle_M G$.
	As $G \subseteq S$ and $\langle E \rangle_M \subseteq S$,
	and so $\langle E \rangle_M G \subseteq S$,
	we conclude $S = \langle E \rangle_M G$.	
\end{proof} 

As $t_1 \cdots t_n g^{-1} = g^{-1} t_1^g \cdots t_n^g$ for a permutation  $g$ and transformations $t_1, \ldots, t_n$  we have,
in \cref{sync_max:prop:factorization_KG} and \cref{sync_max:cor:factorization_EG},
\[ 
\langle K \rangle_M G = G \langle K \rangle_M \text{ and }
\langle E \rangle_M G = G \langle E \rangle_M.
\]
Further, as $G$ contains the identity transformation, we have
$\langle G \cup \{t\} \rangle_S = \langle G \cup \{t\} \rangle_M$.
However, for a non-permutation $t$, $\langle K \rangle_S G \ne \langle K \rangle_M G$,
as $\operatorname{id}_{[n]} \notin \langle K \rangle_S$.

\section{Graphs and Permutation Groups}
\label{prelim:sec:graphs-groups}

%!TEX root = ./../../springer-mono-format/book.tex
%!TEX root = ./../../thesis-trier.tex

Permutation group theory enters graph theory via the automorphism group 
in a natural way. Here, we introduce some further notions
connecting graphs and permutation groups. \Cref{prelim:sec:genSetSn} collects results on generating sets for $\mathcal S_n$. For example, we give a graph-theoretical
characterization when a set of transpositions generate~$\mathcal S_n$.
\Cref{prelim:sec:vertex-transitive}
introduces vertex-transitive graphs and investigates the connectivity
of vertex-transitive graphs. 
%\Cref{prelim:thm:degree-implies-connectivity} \todo{finite?}  
%links the outdegree of a strongly connected vertex-transitive
%graph to its strong $k$-connectivity. 
Lastly, \cref{prelim:sec:orbitalgraphs}
introduces orbital graphs.

We often denote graphs and groups
with the capital letter~$G$, but the context
should make it clear whether we mean a graph or a group.


\subsection{Generating Sets for the Full Symmetric Group}
\label{prelim:sec:genSetSn}

It is a well-known result that permutation cyclically permuting a finite set
and a transposition on this set generate the full symmetric group.
Isaacs and Zieschang~\cite[Thm.~B]{IsaacsZ95} gave the following generalization.

\begin{theorem}[{Isaacs and Zieschang~\cite[Thm.~B]{IsaacsZ95}}]
	\label{prelim:thm:cycles_gen_Sn}
	Let $x = (0 ~ 1 ~ 2 ~ \cdots ~ n-1)_{[n]}$ and $y = (0 ~ 1 ~ 2 ~ \cdots ~ m-1)_{[n]}$ for some $1 < m < n$.
	Then $\langle \{x,y\} \rangle_G$ equals the full symmetric group on $[n]$ unless both $n$ and $m$ are odd, in which
	case $\langle \{x,y\} \rangle_G$ equals $\mathcal A_n$. % is the alternating group.
\end{theorem}

Isaacs and Zieschang~\cite[Thm.~A]{IsaacsZ95} also proved that for $n \ne 4$
and $x \in \mathcal S_n$ there exists $y \in \mathcal S_n$ such that $\langle \{x,y\} \rangle_G = \mathcal S_n$  (the exception being any element $x$ in the Klein four-group from \cref{prelim:permgrp:ex:Kleinfourgroup}).
Actually, this result was first proved by Sophie Piccard~\cite[pp.~80--86]{piccard46} in 1946, 
but got lost in limbo\footnote{We were not able to obtain
	a copy of Piccards original work~\cite{piccard46}
	but learned this through the work of Salomaa~\cite{DBLP:journals/jucs/Salomaa02}.}
and was reproved by Isaacs \& Zieschang~\cite{IsaacsZ95} in 1996.

By \cref{prelim:fact:transpositions_generate_Sn} the full symmetric group on a finite set
is generated by all transpositions. Next, we state a refined result, needed in \cref{sync_max:sec:perm_grps},
that characterizes the sets of transpositions generating $\mathcal S_n$ in terms of an associated graph. 
For the convenience of the reader, we give a proof (which can also be found in~\cite[Lem.~3.10.1]{godsil01};
compare also the proof of~\cite[Lem.~6(a)]{AraujoCameron2014}
which uses a similar argument).

\begin{proposition}
	\label{prop:connected_transpositions_gen_Sn}
	Let $T \subseteq \mathcal S_n$ be a set of transpositions
	and consider the undirected graph $G = ([n], \{ \{a,b\} \mid (a ~ b) \in T\})$.
	Then $G$ is connected if and only if $\langle T \rangle_G = \mathcal S_n$.
\end{proposition}
\begin{proof}
	Define a relation $\sim$ on $[n]$
	by $a \sim b$ if and only if $a = b$ or $(a ~ b) \in \langle T \rangle_G$.
	Clearly, $\sim$ is reflexive and symmetric.
	Further, for $a \sim b$ and $b \sim c$
	with distinct $a,b,c \in [n]$ it is $(a ~ b)^{(b ~ c)} = (a ~ c) \in \langle T \rangle_G$;
	thus $\sim$ is transitive and an equivalence relation.
	As the graph $G$ results from the digraph $G_{\sim}$ by forgetting
	about the direction of the edges, from \cref{prelim:fact:relation_graph}
	we obtain that the connected components equal the equivalence classes of~$\sim$.	
	So, if $G$ is connected, then all points are pairwise comparable for $\sim$, and
	$\langle T \rangle_G$ contains all transpositions.
	By Fact~\ref{prelim:fact:transpositions_generate_Sn}, we find $\langle T \rangle_G = \mathcal S_n$.
	Otherwise, suppose~$G$ is not connected.
	Then, we can pick $a,b$
	with $a \not\sim b$, i.e., $(a ~ b) \notin \langle T \rangle_G$
	and so~$\langle T \rangle_G \ne \mathcal S_n$.
\end{proof}

\begin{remark}	
	In the proof of Jordan's theorem from~\cite{IsaacsZ95},
	it is also shown that for a set of $3$-cycles $T$, $\langle T \rangle_G$ contains the alternating
	group if and only if the graph on~$[n]$
	where $a,b \in [n]$ are joined by an edge if and only if 
	$a,b \in \move{t}$ for some $t \in T$
	is connected.
\end{remark}

In~\cite{Knoerr2007} it was shown that a $k$-transitive group of degree $n > k$
is generated by all elements fixing precisely~$i$ points if $i \le k - 3$.
Since the full symmetric group on a set of size~$n$ is $k$-transitive for each $1 \le k \le n$,
by this result we obtain more generating sets for the full symmetric group of a finite set.

\subsection{Vertex-Transitive Graphs}
\label{prelim:sec:vertex-transitive}

\IN{vertex-transitive graph}
\IN{outdegree of graph}
\IN{strong connectivity $\kappa(G)$}
\label{def:kappa}
A (directed or undirected) graph is \emph{vertex-transitive} if its automorphism group acts transitively on its
vertices.
If~$G = (V,E)$ is vertex-transitive, then every vertex has the same outdegree 
and we call the common value the \emph{outdegree of~$G$}.
%The \emph{strong connectivity}~$\kappa(G)$ is the unique number~$k$ \todo{warum $K_n$ ausgeklammert? weil dann nie zwei SCCs}
%such that % K_n sagen einfach teil der ausagen.
%$G$ is strongly~$k$-connected but not strongly $(k+1)$-connected.

A proof of the next statement can be found in~\cite[Thm~1.10]{cameron_1999},
\cite[Thm~1.1]{ArCamSt2017}, \cite[Chap.~2.6]{godsil01},
and in~\cite{Rystsov2022} (in the latter work formulated with slightly different notions but easily seen to give the result).

\begin{theorem}
	\label{prelim:thm:vertex_transitive_finite_graph}
	Let $G = (V, E)$ be a finite vertex-transitive digraph.
	If $G$ is weakly connected, then $G$ is strongly connected.
\end{theorem}

Next, we link the outdegree of a vertex-transitive digraph to its connectivity. The result appears in~\cite[Prop.~3.4]{Hamidoune81} (compare also~\cite[Thm~3.7 (a)]{BabaiSurvey96} where it is formulated
in terms of the connectivity 
from \cref{prelim:rem:connectedconnectivity}).

\begin{theorem}[{Hamidoune~\cite{Hamidoune81}}]
	\label{prelim:thm:degree-implies-connectivity}
	If $G = (V, E)$
	is a finite strongly connected vertex-transitive digraph of outdegree~$d$ without edges of the form $(v, v)$ for $v \in V$,
	then~$G$ is strongly
	$\left( \left\lfloor \frac{d}{2} \right\rfloor + 1 \right)$-connected.
\end{theorem} 
\begin{proof} \label{def:Nplus} \IN{atom in a graph}
	We can assume $E \ne V \times V$. 
	Since~$G$ is finite, there exists a largest number~$k$ 
	such that~$G$ is 
	strongly $k$-connected, but not strongly $(k+1)$-connected.
	For $U \subseteq V$
	set $N^+(U) = \{ v \in V \setminus U \mid \exists u \in U : (u,v) \in E \}$. 
	We have $k \le |N^+(\{u\})| = d$ for each $u \in V$. % A = N^+(u) muss jeder pfad zu nem anderen schneiden, also nicht ( |N^+(u)|+1 )-connected, damit | | \le kappe
	An \emph{atom}~\cite{hamidoune77} $A \subseteq V$ is a set of minimal cardinality with
	$|N^+(A)| = k$ and
	$A \cup N^+(A) \ne V$.\footnote{For undirected graphs, atoms 
		were introduced in~\cite{Mader70,Mader71,Watkins70}.}
	In a strongly connected digraph the atoms partition~$V$~\cite{hamidoune77}
	(see also~\cite[Thm. 2.1 in both]{DBLP:journals/ejc/Hamidoune81,Hamidoune81}).
	As an automorphism maps atoms to atoms,
	the atoms induce an $\operatorname{Aut}(G)$-congruence on the vertices.
	Hamidoune proved that $|A| \le k$~\cite[Lem.~3.3]{Hamidoune81}.
	So, for $u \in A$, as $N^+(\{u\}) \subseteq A \setminus \{u\} \cup N^+(A)$, we get 
	%	$ 
	%	d \le |A \setminus \{u\}| + |N^+(\{u\})| < |A| + |N^+(A)| = |A| + \kappa \le 2\kappa. 
	%	$
	% N^+ nicht monoton bzgl submengen (beispiel dreieck mit allen richtungen und 2-er menge und 3-er menge)
	$
	d \le |A \setminus \{u\}| + |N^+(A)| < |A| + k \le 2k.
	$ % u betrachten, geht entweder raus, oder bleibt in A zu den A \ {u} anderen, damit folgt erste Ungleichung.
\end{proof} 

In the context of the present work, we actually only need the following
implication.

\begin{corollary}
	\label{prelim:prop:vertex-transitive-SCC-32-connected}
	If $G = (V, E)$
	is a finite strongly connected vertex-transitive graph without edges of the form $(v,v)$ for $v \in V$,
	then $G$ is strongly $\frac{3}{2}$-connected.
\end{corollary}
\begin{proof}
	Let $d$ be the outdegree of $G$.	
	For $d = 1$, $G$ forms a single cycle,
	i.e., a strongly connected digraph
	where every vertex has outdegree one,	
	which is clearly strongly~$\frac{3}{2}$-connected.
	For $d \ge 2$,
	by \cref{prelim:thm:degree-implies-connectivity}, 
	$G$ is strongly~$2$-connected, hence strongly $\frac{3}{2}$-connected.	
\end{proof} 

Note that, by \cref{prelim:prop:vertex-transitive-SCC-32-connected},
the conclusion of \cref{prelim:thm:degree-implies-connectivity}
is valid for $d = 1$ even after removal of the 
rounding, which somehow explains
our choice of naming for strong $\frac{3}{2}$-connectivity.

\subsection{Orbital Graphs}
\label{prelim:sec:orbitalgraphs}

Suppose $G \le_G \mathcal S_X$ is a permutation group.
Let $a,b \in X$. Then we can associate two graphs with the tuple $(a,b)$
in a natural way.

\IN{orbital graph!directed}
\IN{orbital graph!undirected}
\begin{enumerate} 
	\item The \emph{directed orbital graph} of the tuple $(a,b)$
	is the graph with vertex set~$X$ and edge set $\{ (\applyfunc{a}{g}, \applyfunc{b}{g}) \mid g \in G \}$.
	
	\item The \emph{undirected orbital graph} of the tuple $(a,b)$ (or the $2$-set $\{a,b\}$)
	is the graph with vertex set~$X$ and edge set $\{ \{ \applyfunc{a}{g}, \applyfunc{b}{g} \} \mid g \in G \}$.	
\end{enumerate} 

\begin{example} 
 \Cref{fig:orbitals} shows the directed orbital graphs
 of the group $G = \langle \{ (0 ~ 1 ~ 2 ~ 3)\} \rangle_G$.
 The undirected orbital graphs are obtained from these
 by forgetting about the directions on the edges.
 Since~$G$ contains no element swapping two points, 
 the orbital graphs are oriented.

\todo[inline]{graphiken zu groß}

\begin{figure}[th]
	\centering
	\begin{tikzpicture}[>=Stealth, node distance=2em, every node/.style={minimum size=0.5cm}]
		\node (1) at (0,2) {0};
		\node (2) at (2,2) {1};
		\node (3) at (2,0) {2};
		\node (4) at (0,0) {3};
		
		\path[->] (1) edge [loop above] (1);
		\path[->] (2) edge [loop above] (2);
		\path[->] (3) edge [loop above] (3);
		\path[->] (4) edge [loop above] (4);
		
		\node (1r) at (4,2) {0};
		\node (2r) at (6,2) {1};
		\node (3r) at (6,0) {2};
		\node (4r) at (4,0) {3};
		
		\draw (1r) -- (2r); \draw (2r) -- (3r); \draw (3r) -- (4r); \draw (4r) -- (1r);
				
		\node (1s) at (8,2) {0};
		\node (2s) at (10,2) {1};
		\node (3s) at (10,0) {2};
		\node (4s) at (8,0) {3};
		
		\draw (1s) -- (3s); \draw (4s) -- (2s);
		\draw (3s) -- (1s); \draw (2s) -- (4s);
		
		\node (1t) at (12,2) {0};
		\node (2t) at (14,2) {1};
		\node (3t) at (14,0) {2};
		\node (4t) at (12,0) {3};
		
		\draw (1t) -- (4t); \draw (4t) -- (3t); \draw (3t) -- (2t); \draw (2t) -- (1t);
		
	\end{tikzpicture}
	\caption[Directed orbital graphs of $G = \langle \{ (0 ~ 1 ~ 2 ~ 3)\} \rangle_G$]{The directed orbital graphs of $G = \langle \{ (0 ~ 1 ~ 2 ~ 3)\} \rangle_G$.}
	\label{fig:orbitals}
\end{figure}
\end{example}

\begin{example}
	\Cref{fig:orbitals-klein-four} shows the directed orbital graphs
	of the group $G = \langle \{ (0 ~ 1)(2 ~ 3), (0 ~ 2)(1 ~ 3)\} \rangle_G$.
	All of them except the left-most one are symmetric.
	
	\begin{figure}[th]
		\centering
		\begin{tikzpicture}[>=triangle 45, node distance=1.5cm, every node/.style={minimum size=0.5cm}]
			
			\node (1) at (0,2) {0};
			\node (2) at (2,2) {1};
			\node (3) at (2,0) {2};
			\node (4) at (0,0) {3};
			
			\path[->] (1) edge [loop above] (1);
			\path[->] (2) edge [loop above] (2);
			\path[->] (3) edge [loop above] (3);
			\path[->] (4) edge [loop above] (4);
			
			\node (1r) at (4,2) {0};
			\node (2r) at (6,2) {1};
			\node (3r) at (6,0) {2};
			\node (4r) at (4,0) {3};
			
			\draw[<->] (1r) -- (2r);
			\draw[<->] (3r) -- (4r);
			
			\node (1s) at (8,2) {0};
			\node (2s) at (10,2) {1};
			\node (3s) at (10,0) {2};
			\node (4s) at (8,0) {3};
			
			\draw[<->] (1s) -- (3s);
			\draw[<->] (2s) -- (4s);
						
			\node (1t) at (12,2) {0};
			\node (2t) at (14,2) {1};
			\node (3t) at (14,0) {2};
			\node (4t) at (12,0) {3};
			
			\draw[<->] (1t) -- (4t);
			\draw[<->] (2t) -- (3t);
			
		\end{tikzpicture}
		\caption[Directed orbital graphs of $G = \langle \{ (0 ~ 1)(2 ~ 3), (0 ~ 2)(1 ~ 3)\} \rangle_G$]{The directed orbital graphs of $G = \langle \{ (0 ~ 1)(2 ~ 3), (0 ~ 2)(1 ~ 3)\} \rangle_G$.}
		\label{fig:orbitals-klein-four}
	\end{figure}
\end{example}

The edge set of the directed orbital graph of $(a,b)$
is contained in $\Orbit{a}{G} \times \Orbit{b}{G}$.
The directed orbital graph of $(a,b)$ 
is either symmetric
or oriented, and the former case applies precisely if $G$
contains a permutation swapping~$a$ and~$b$~\cite[Lem.~2.4.2]{godsil01}.
Also, $G$ is intransitive if and only if 
in the (directed or undirected) orbital graph
exist vertices that are not part of any edge.
Furthermore,~$G$ is a subgroup of the automorphism
group of each (directed or undirected) orbital graph.
So, by Theorem~\ref{prelim:thm:vertex_transitive_finite_graph} 
the weakly connected components of a finite directed orbital graph are precisely
the strongly connected components.
Finally, a (directed or undirected)
graph is vertex-transitive
if and only if it is a union of orbital graphs w.r.t.~its automorphism group.

\begin{lemma} \label{prelim:lem:orbitals-distinct-orbits}
	Let $G \le_G \mathcal S_n$ and $a,b \in [n]$.
	Then $\Orbit{a}{G} \ne \Orbit{b}{G}$
	if and only if in the undirected orbital graph of $\{a,b\}$ (resp.~directed orbital graph of $(a,b)$)
	there exists no edge with 
	both 
	endpoints in $\Orbit{a}{G}$
	or both endpoints in $\Orbit{b}{G}$
	(resp.~the starting vertex of each edge is in $\Orbit{a}{G}$ 
	and its ending vertex is in $\Orbit{b}{G}$).
\end{lemma}
\begin{proof}
	Since $\Orbit{a}{G}$ and $\Orbit{b}{G}$ are orbits, we have $\Orbit{a}{G} = \Orbit{b}{G}$
	if and only if there exists $g \in G$ such that $g(a) \in \Orbit{b}{G}$ or $g(b) \in \Orbit{a}{G}$.
	So,  $\{g(a), g(b)\} \subseteq \Orbit{a}{G}$
	or $\{g(a), g(b)\} \subseteq \Orbit{b}{G}$ for some $g \in G$
	if and only if~$G$ is transitive, which
	gives the claim for the undirected orbital graph.
	Furthermore, since the edge set of the directed orbital graph is in $\Orbit{a}{G} \times \Orbit{b}{G}$,
	if $G$ is intransitive every edge of the directed orbital graph
	starts in $\Orbit{a}{G}$ and ends in $\Orbit{b}{G}$.	
\end{proof}

\section{Discussion}
\label{prelim:permgrp:sec:discussion}

\todo[inline]{diskussion in datei und oben kommentierten teil raus}
Jordan~\cite{jordan1870} wrote the first textbook on permutation group theory (but with an antique notation and terminology).
Prior to this, permutation groups in the context of polynomial
equations (as permutations of their root)
were investigated by Lagrange, Galois, Abel and Ruffini, see~\cite{Kleiner2007_HistoryAbstractAlgebra,Neumann2011,Wussing1984_GenesisAbstractGroup}.
At the end of the 19th century, Burnside published his classical textbook~\cite{burnside1911} (1st edition 1897) 
which collected and contained most results from 19th century group theory with a focus on permutation groups.
Wielandt~\cite{wielandt1964finite} wrote
a standard textbook on permutation group theory
which established some standard terminology.
Other standard introductions are the textbooks by Dixon \& Mortimer~\cite{dixon1996permutation}
and by Cameron~\cite{cameron_1999}. 
Interestingly, permutation group theory, in its historical development, preceded semigroup theory, 
see~\cite{Hollings2009}.

Before the middle of the 19th century, 
groups were synonymous to permutation groups on a finite set
and were simply thought of as non-empty subsets of permutations closed under function composition (called substitution groups back then). The concept of an abstract group slowly emerged.
Cayley~\cite{Cayley54a,Cayley54b,Cayley59} was the first who
viewed groups more abstract solely in terms of laws, and not in some permutation representation.
However, his conceptualization was still far away from our modern view and still relied
on implicit assumptions. Weber gave an axiom system
for finite groups~\cite{Weber1882_MathAnn} and for general groups~\cite{Weber1896_LehrbuchAlgebra_Vol2}. 
During the 19th century many further famous mathematician besides the ones already mentioned 
were involved into forming
the abstract theory of groups as, for example, Cauchy, Kronecker, Dyck, Frobenius, and Hölder
to name a few. The first definition of an abstract group
in its modern form appears in van der Waerden's textbook~\cite{vanderWaerden_ModernAlgebra_SpringerArchive}, which is based
on lectures by Artin and Noether from the beginning of the 20th century.
For more information on the history and genesis of abstract group theory, see~\cite{Kleiner2007_HistoryAbstractAlgebra,Neumann_AbstractGroups_MacTutor,Wussing1984_GenesisAbstractGroup}.

Not all authors use the term $G$-congruence (some just call them equivalence relations preserved by the group), 
but it is a notion that appears in the literature, for example in~\cite{Burness14}.
The generalization of it to transformation semigroups is obvious, but I have not
seen in the literature. In~\cite{KilpKnauerMikhalev2000_MonoidsActs}, in the more general context
of semigroup actions, the authors introduce the term $S$-act congruences,
and $G$-congruences are also used in the more general context of group actions, also called $G$-sets,
for example, in~\cite{NeumannStoyThompson1994GroupsAndGeometry}.

The books~\cite{Seress2003PermutationGroupAlgorithms,Butler1991AlgorithmsPermutationGroups,Holt2005_ComputationalGroupTheory} contains information on algorithms for groups and permutation groups. Beaudry~\cite{Beaudry1989CharacterizationIdempotentTM} 
investigated the computational complexity of deciding whether all transformations in a 
given transformation monoid are idempotent.

\IN{point-symmetric graph}
Vertex-transitive graphs seemed to be introduced
by Sabidussi~\cite{sabidussi1964vertex}, at least this is the first reference
that we were able to find where this term was introduced.
Vertex-transitive graphs are also called \emph{point-symmetric} 
in~\cite{harary1969graph}.
The study of symmetry in graphs was initiated 
by Foster~\cite{foster1932geometrical}.
More information on vertex-transitive graphs
can be found in~\cite{BabaiSurvey96,Biggs1993,Cameron2004_AutomorphismsOfGraphs,PraegerLiNiemeyer1997_FiniteTransitivePermutationGroups}.

